\documentclass[12pt,a4paper]{article}
\usepackage{fontspec}
\usepackage{geometry}
\usepackage{enumitem}
\usepackage{fancyhdr}
\usepackage{lastpage}
\usepackage{hyperref}

% Configure IBM Plex Sans font family
\setmainfont{IBMPlexSans}[
  Path=/home/asi0/asi0-repos/bitcoin-educational-content/scripts/quizz_to_pdfs/fonts/TrueType/IBM-Plex-Sans/,
  Extension=.ttf,
  UprightFont=*-Regular,
  BoldFont=*-Bold,
  Scale=1.0
]

\geometry{margin=2.5cm}
\pagestyle{fancy}
\fancyhf{}
\rhead{Page \thepage\ of \pageref{LastPage}}
\lhead{HIS203 - Quiz (fr)}
\renewcommand{\headrulewidth}{0.4pt}

\title{Quiz: HIS203}
\author{Language: fr}
\date{\today}

\begin{document}

\maketitle
\thispagestyle{fancy}

\section*{Instructions}
Answer all questions by selecting the best option (A, B, C, or D). The answer key is provided at the end of this document.

\bigskip


\subsection*{Question 1}
Quel pseudonyme Ross Ulbricht utilisait-il pour faire la promotion de Silk Road sur les forums en ligne ?

\begin{enumerate}[label=\Alph*., leftmargin=*]
\item Teppy
\item ShadowofHarbringer
\item Dread Pirate Roberts
\item Altoid
\end{enumerate}

\bigskip


\subsection*{Question 2}
Quelles sont les deux innovations techniques principales qui ont permis à Silk Road de se démarquer des autres places de marché illégales existantes ?

\begin{enumerate}[label=\Alph*., leftmargin=*]
\item Pecunix et Liberty Reserve
\item Tor et Bitcoin
\item eBay et Amazon
\item PayPal et Western Union
\end{enumerate}

\bigskip


\subsection*{Question 3}
Quel était le nom du blog Wordpress que Ross Ulbricht avait créé pour donner des instructions d'accès à Silk Road ?

\begin{enumerate}[label=\Alph*., leftmargin=*]
\item silkroadguide.wordpress.com
\item silkroad2011.wordpress.com
\item silkroadmarket.wordpress.com
\item silkroad420.wordpress.com
\end{enumerate}

\bigskip


\subsection*{Question 4}
Quel était le nom de l'entreprise de revente de livres d'occasion que Ross Ulbricht gérait en 2010 ?

\begin{enumerate}[label=\Alph*., leftmargin=*]
\item Underground Books
\item Digital Library Store
\item Good Wagon Books
\item Austin Book Exchange
\end{enumerate}

\bigskip


\subsection*{Question 5}
Quel philosophe a théorisé l'agorisme dans les années 1970, une doctrine qui a profondément influencé Ross Ulbricht ?

\begin{enumerate}[label=\Alph*., leftmargin=*]
\item Murray Rothbard
\item Samuel Edward Konkin III
\item Joseph Salerno
\item Ludwig von Mises
\end{enumerate}

\bigskip


\subsection*{Question 6}
Quel était le nom du réseau superposé confidentiel utilisé par Silk Road pour permettre la navigation anonyme ?

\begin{enumerate}[label=\Alph*., leftmargin=*]
\item I2P
\item Tor
\item VPN
\item Freenet
\end{enumerate}

\bigskip


\subsection*{Question 7}
Quel était le surnom donné à la plateforme Silk Road en raison de sa popularité dans le commerce en ligne de produits illicites ?

\begin{enumerate}[label=\Alph*., leftmargin=*]
\item L'eBay du darknet
\item Le PayPal noir
\item L'Amazon de la drogue
\item Le Walmart souterrain
\end{enumerate}

\bigskip


\subsection*{Question 8}
Quel était le nom du système de devise numérique que Ross Ulbricht envisageait initialement d'utiliser pour Silk Road avant de découvrir Bitcoin ?

\begin{enumerate}[label=\Alph*., leftmargin=*]
\item PayPal
\item Liberty Reserve
\item Pecunix
\item E-gold
\end{enumerate}

\bigskip


\subsection*{Question 9}
Quel était le principal moyen d'expédition utilisé par Silk Road pour livrer les produits aux acheteurs ?

\begin{enumerate}[label=\Alph*., leftmargin=*]
\item Le service postal classique (comme l'United States Postal Service)
\item Des services de livraison décentralisés utilisant la blockchain
\item Des coursiers privés spécialisés dans l'anonymat
\item Un réseau de distribution physique via des points de collecte
\end{enumerate}

\bigskip


\subsection*{Question 10}
Quelle était la motivation principale de Ross Ulbricht pour créer Silk Road selon ses propres déclarations ?

\begin{enumerate}[label=\Alph*., leftmargin=*]
\item Générer des profits importants grâce au commerce de drogues illégales
\item Créer une alternative légale aux pharmacies traditionnelles
\item Tester les capacités techniques de Bitcoin et du réseau Tor
\item Réduire l'influence de l'État dans la société grâce à l'économie souterraine
\end{enumerate}

\bigskip


\subsection*{Question 11}
Quelle erreur fondamentale Ross Ulbricht a-t-il commise en octobre 2011 qui l'aidera plus tard à être identifié par les autorités ?

\begin{enumerate}[label=\Alph*., leftmargin=*]
\item Il a révélé l'adresse physique de son serveur dans un message sur le forum
\item Il a donné une interview télévisée en révélant accidentellement son identité
\item Il a utilisé son adresse email contenant son nom civil dans une annonce de recrutement
\item Il a publié une photo de lui-même sur la page d'accueil de Silk Road
\end{enumerate}

\bigskip


\subsection*{Question 12}
Quel était le nom du concurrent de Silk Road qui a émergé à la fin du mois de juin 2011, copiant son modèle ?

\begin{enumerate}[label=\Alph*., leftmargin=*]
\item Underground Brokers
\item Black Market Reloaded
\item Adamflowers
\item Dark Web Market
\end{enumerate}

\bigskip


\subsection*{Question 13}
Qui était Mark Karpelès avant de reprendre Mt. Gox en 2011?

\begin{enumerate}[label=\Alph*., leftmargin=*]
\item Un expert en sécurité informatique qui avait créé plusieurs systèmes de paiement numérique
\item Un développeur français passionné de technologie qui gérait un service d'hébergement au Japon
\item Un banquier spécialisé dans les cryptomonnaies qui travaillait pour une institution financière
\item Un entrepreneur américain expérimenté dans le domaine des plateformes d'échange en ligne
\end{enumerate}

\bigskip


\subsection*{Question 14}
Quel était le montant approximatif des bitcoins volés lors du krach éclair de Mt. Gox le 19 juin 2011?

\begin{enumerate}[label=\Alph*., leftmargin=*]
\item 25 000 BTC (environ 500 000 dollars)
\item 80 000 BTC (environ 1,4 million de dollars)
\item 2 000 BTC (environ 35 000 dollars)
\item 17 000 BTC (environ 235 000 dollars)
\end{enumerate}

\bigskip


\subsection*{Question 15}
Quel était le principal motif d'intérêt de Mark Karpelès pour Bitcoin selon ses propres déclarations?

\begin{enumerate}[label=\Alph*., leftmargin=*]
\item Les motivations idéologiques libertariennes
\item Les aspects techniques et technologiques du système
\item Les opportunités de profit et de spéculation
\item La possibilité de créer une monnaie alternative
\end{enumerate}

\bigskip


\subsection*{Question 16}
Quelle était la première plateforme de change européenne à faire réellement concurrence à Mt. Gox en 2011?

\begin{enumerate}[label=\Alph*., leftmargin=*]
\item Bitcoin-Central, lancée en France par David François
\item Bitomat, lancée en Pologne pour les zlotys
\item Britcoin, lancée au Royaume-Uni par Amir Taaki
\item VirWoX, lancée en Autriche pour les monnaies virtuelles
\end{enumerate}

\bigskip


\subsection*{Question 17}
Combien de bitcoins Mt. Gox avait-elle perdus lors du vol du 1er mars 2011, quelques jours avant l'officialisation du transfert de propriété?

\begin{enumerate}[label=\Alph*., leftmargin=*]
\item 80 000 bitcoins
\item 17 000 bitcoins
\item 50 000 bitcoins
\item 25 000 bitcoins
\end{enumerate}

\bigskip


\subsection*{Question 18}
Quel service Mark Karpelès gérait-il avant de reprendre Mt. Gox?

\begin{enumerate}[label=\Alph*., leftmargin=*]
\item KalyHost, un service d'hébergement web
\item Bitcoin.it, le wiki principal de Bitcoin
\item Tibanne Co. Ltd., une société de développement
\item QBitcoin, une implémentation logicielle alternative
\end{enumerate}

\bigskip


\subsection*{Question 19}
Quel était le prix du bitcoin lors du pic de la première bulle spéculative en juin 2011?

\begin{enumerate}[label=\Alph*., leftmargin=*]
\item 32 dollars
\item 10 dollars
\item 50 dollars
\item 100 dollars
\end{enumerate}

\bigskip


\subsection*{Question 20}
Quel était le nom de la société que Mark Karpelès avait fondée au Japon avant de reprendre Mt. Gox?

\begin{enumerate}[label=\Alph*., leftmargin=*]
\item Tibanne Co. Ltd.
\item Bitcoin Consultancy Ltd.
\item KalyHost Co. Ltd.
\item QBitcoin Co. Ltd.
\end{enumerate}

\bigskip


\subsection*{Question 21}
Quel était le nom de la plateforme de change britannique lancée par Amir Taaki en mars 2011?

\begin{enumerate}[label=\Alph*., leftmargin=*]
\item Intersango
\item Britcoin
\item TradeHill
\item Bitcoin-Central
\end{enumerate}

\bigskip


\subsection*{Question 22}
Combien d'utilisateurs Mt. Gox comptait-elle approximativement à la fin du mois de mai 2011?

\begin{enumerate}[label=\Alph*., leftmargin=*]
\item 25 000 utilisateurs
\item 60 000 utilisateurs
\item 100 000 utilisateurs
\item 35 000 utilisateurs
\end{enumerate}

\bigskip


\subsection*{Question 23}
Quel était le nom du documentaire français de 2007 dans lequel Mark Karpelès est apparu avant de devenir célèbre avec Mt. Gox?

\begin{enumerate}[label=\Alph*., leftmargin=*]
\item Tech Revolution
\item Geek Generation
\item Digital Dreams
\item Suck my Geek
\end{enumerate}

\bigskip


\subsection*{Question 24}
Quel était le nom du service de dépôt fiduciaire géré par Gavin Andresen qui a été utilisé pour placer en séquestre les bitcoins achetés lors du krach éclair de Mt. Gox?

\begin{enumerate}[label=\Alph*., leftmargin=*]
\item SafeDeposit
\item BitcoinEscrow
\item TrustCoin
\item ClearCoin
\end{enumerate}

\bigskip


\subsection*{Question 25}
Quel surnom a été donné à Roger Ver en raison de son promotion zélée de Bitcoin?

\begin{enumerate}[label=\Alph*., leftmargin=*]
\item Le Prophète de Bitcoin
\item Le Jésus de Bitcoin
\item L'Évangéliste de Bitcoin
\item Le Messie de Bitcoin
\end{enumerate}

\bigskip


\subsection*{Question 26}
Quel mouvement politique américain a particulièrement bénéficié à l'adoption de Bitcoin selon le chapitre?

\begin{enumerate}[label=\Alph*., leftmargin=*]
\item Le mouvement démocrate progressiste
\item Le mouvement libertarien américain
\item Le mouvement républicain conservateur
\item Le mouvement socialiste américain
\end{enumerate}

\bigskip


\subsection*{Question 27}
Quelle émission de radio libertarienne a joué un rôle clé dans la découverte de Bitcoin par Roger Ver?

\begin{enumerate}[label=\Alph*., leftmargin=*]
\item InfoWars
\item The Bitcoin Show
\item Free Talk Live
\item Liberty Radio Network
\end{enumerate}

\bigskip


\subsection*{Question 28}
Quel État américain a été choisi par le Free State Project pour regrouper les libertariens?

\begin{enumerate}[label=\Alph*., leftmargin=*]
\item Le New Hampshire
\item Le Texas
\item Le Colorado
\item La Californie
\end{enumerate}

\bigskip


\subsection*{Question 29}
Quelle école de pensée économique est particulièrement appréciée par les libertariens américains?

\begin{enumerate}[label=\Alph*., leftmargin=*]
\item L'école monétariste
\item L'école de Chicago
\item L'école autrichienne
\item L'école keynésienne
\end{enumerate}

\bigskip


\subsection*{Question 30}
Quel slogan caractérisait le mouvement Occupy Wall Street qui a utilisé Bitcoin pour contourner la censure financière?

\begin{enumerate}[label=\Alph*., leftmargin=*]
\item Nous sommes les 99 %
\item Occupons la Réserve fédérale
\item Vivre libre ou mourir
\item La séparation de la monnaie et de l'État
\end{enumerate}

\bigskip


\subsection*{Question 31}
Pourquoi certains libertariens traditionnels ont-ils initialement rejeté Bitcoin malgré leurs convictions anti-étatiques?

\begin{enumerate}[label=\Alph*., leftmargin=*]
\item Parce qu'il était contrôlé par des entreprises privées plutôt que par le marché
\item Parce qu'il nécessitait une connexion internet permanente pour fonctionner
\item Parce qu'il était trop complexe techniquement pour être adopté massivement
\item Parce qu'il n'avait pas de valeur intrinsèque selon le théorème de régression de Ludwig von Mises
\end{enumerate}

\bigskip


\subsection*{Question 32}
Quel évènement a motivé la création du mouvement Occupy Wall Street qui a ensuite adopté Bitcoin?

\begin{enumerate}[label=\Alph*., leftmargin=*]
\item La campagne présidentielle de Ron Paul en 2012
\item La fermeture du site Silk Road par les autorités
\item L'effondrement de la plateforme d'échange Mt. Gox
\item La crise des subprimes et la récession économique qui a suivi
\end{enumerate}

\bigskip


\subsection*{Question 33}
Quelle plateforme publicitaire ont lancé Erik Voorhees et Ira Miller en novembre 2011 pour promouvoir du contenu sur Twitter en utilisant Bitcoin?

\begin{enumerate}[label=\Alph*., leftmargin=*]
\item Coinapult
\item BitMunchies
\item BTCinch
\item FeedZeBirds
\end{enumerate}

\bigskip


\subsection*{Question 34}
Quel était le surnom de Ron Paul en raison de sa tendance à voter contre les projets de loi qui réduisent la liberté ou augmentent les dépenses publiques?

\begin{enumerate}[label=\Alph*., leftmargin=*]
\item Mister Freedom
\item Doctor No
\item Doctor Liberty
\item The Libertarian
\end{enumerate}

\bigskip


\subsection*{Question 35}
Quelle était la devise officielle de l'État du New Hampshire qui a influencé le choix du Free State Project?

\begin{enumerate}[label=\Alph*., leftmargin=*]
\item Freedom First (La liberté d'abord)
\item Liberty or Death (La liberté ou la mort)
\item Don't Tread on Me (Ne me marchez pas dessus)
\item Live Free or Die (Vivre libre ou mourir)
\end{enumerate}

\bigskip


\subsection*{Question 36}
Quel était le principal reproche des libertariens traditionnels envers Bitcoin selon le théorème de régression de Ludwig von Mises?

\begin{enumerate}[label=\Alph*., leftmargin=*]
\item Bitcoin n'avait pas de valeur intrinsèque ni d'adossement à quelque chose de concret
\item Bitcoin nécessitait l'intervention de l'État pour fonctionner correctement
\item Bitcoin était contrôlé par une autorité centrale gouvernementale
\item Bitcoin permettait une inflation monétaire incontrôlée par les banques centrales
\end{enumerate}

\bigskip


\subsection*{Question 37}
Quel événement a déclenché le premier véritable marché baissier du Bitcoin en 2011?

\begin{enumerate}[label=\Alph*., leftmargin=*]
\item La publication de l'article critique de Paul Krugman dans le New York Times
\item L'interdiction du trading de Bitcoin par la Securities and Exchange Commission
\item La faillite de la plateforme de change TradeHill et l'arrêt des services Dwolla
\item Le piratage de Mt. Gox en juin suivi de la fermeture de MyBitcoin en juillet
\end{enumerate}

\bigskip


\subsection*{Question 38}
Quel était le prix du Bitcoin en novembre 2011 après la chute de plus de 90% par rapport au sommet de juin?

\begin{enumerate}[label=\Alph*., leftmargin=*]
\item Environ 6 dollars
\item Environ 2 dollars
\item Environ 15 dollars
\item Environ 10 dollars
\end{enumerate}

\bigskip


\subsection*{Question 39}
Quel était le nom du service de change instantané créé par Gareth Nelson et Charlie Shrem pour faciliter les dépôts vers les plateformes d'échange?

\begin{enumerate}[label=\Alph*., leftmargin=*]
\item InstantCoin
\item BitInstant
\item QuickBitcoin
\item BitExchange
\end{enumerate}

\bigskip


\subsection*{Question 40}
Qui était le jeune développeur chinois de 17 ans qui a créé la plateforme de trading sur marge Bitcoinica en septembre 2011?

\begin{enumerate}[label=\Alph*., leftmargin=*]
\item Zhou Tong (Ryan Tong Zhou)
\item Mircea Popescu
\item James McCarthy (Nefario)
\item Charlie Shrem
\end{enumerate}

\bigskip


\subsection*{Question 41}
Quel était le nom de la plateforme de change créée par James McCarthy (Nefario) qui permettait d'émettre des actions et des obligations en bitcoins?

\begin{enumerate}[label=\Alph*., leftmargin=*]
\item Bitcoinica
\item MPEx (Mircea Popescu's Exchange)
\item BitInstant
\item GLBSE (GLobal Bitcoin Stock Exchange)
\end{enumerate}

\bigskip


\subsection*{Question 42}
Quel était le prix approximatif du Bitcoin au sommet de la bulle spéculative de juin 2011?

\begin{enumerate}[label=\Alph*., leftmargin=*]
\item 64 dollars
\item 48 dollars
\item 32 dollars
\item 16 dollars
\end{enumerate}

\bigskip


\subsection*{Question 43}
Quelle était la principale difficulté que BitInstant tentait de résoudre pour les utilisateurs occidentaux de Mt. Gox?

\begin{enumerate}[label=\Alph*., leftmargin=*]
\item Les utilisateurs ne pouvaient pas créer de comptes depuis l'Occident
\item Les frais de transaction étaient trop élevés sur la plateforme japonaise
\item La plateforme Mt. Gox ne supportait pas les devises occidentales
\item Les virements internationaux vers le Japon prenaient plusieurs jours d'attente
\end{enumerate}

\bigskip


\subsection*{Question 44}
Quelle expression est devenue populaire dans la communauté Bitcoin pour décrire le fait de perdre de l'argent par liquidation sur la plateforme Bitcoinica?

\begin{enumerate}[label=\Alph*., leftmargin=*]
\item Se faire leverager (get Leveraged)
\item Se faire zhou-tonguer (get Zhou Tonged)
\item Se faire bitcoiniquer (get Bitcoinica'd)
\item Se faire mt-goxer (get Mt Goxed)
\end{enumerate}

\bigskip


\subsection*{Question 45}
Quel était l'objectif principal de la création de plateformes comme GLBSE et MPEx dans l'écosystème Bitcoin de 2011-2012?

\begin{enumerate}[label=\Alph*., leftmargin=*]
\item Permettre l'émission d'actions et d'obligations pour contourner les réglementations financières traditionnelles
\item Créer des portefeuilles sécurisés pour stocker les bitcoins des utilisateurs particuliers
\item Faciliter les échanges de bitcoins contre des monnaies fiduciaires avec moins de frais
\item Offrir des services de trading sur marge avec effet de levier pour les spéculateurs
\end{enumerate}

\bigskip


\subsection*{Question 46}
Quelle était la principale innovation financière apportée par la plateforme Bitcoinica créée par Zhou Tong en septembre 2011?

\begin{enumerate}[label=\Alph*., leftmargin=*]
\item Le transfert rapide de fonds vers Mt. Gox sans virements internationaux
\item L'émission d'actions et d'obligations pour les entreprises de l'écosystème Bitcoin
\item Le trading sur marge avec effet de levier et la possibilité de vente à découvert
\item L'échange instantané de bitcoins contre des dollars sans frais de transaction
\end{enumerate}

\bigskip


\subsection*{Question 47}
Quel était l'argument principal de Paul Krugman contre Bitcoin dans son article 'Golden Cyberfetters' publié en septembre 2011?

\begin{enumerate}[label=\Alph*., leftmargin=*]
\item Bitcoin manquait de soutien gouvernemental et de garanties institutionnelles pour être viable à long terme
\item Bitcoin créait un système déflationniste similaire à l'étalon-or qui encourageait la thésaurisation au lieu des échanges
\item Bitcoin était trop volatile pour servir de réserve de valeur stable dans l'économie mondiale
\item Bitcoin consommait trop d'énergie pour être une solution monétaire durable et écologiquement responsable
\end{enumerate}

\bigskip


\subsection*{Question 48}
Quelle méthode de financement Mircea Popescu a-t-il utilisée en février 2012 pour développer sa plateforme MPEx?

\begin{enumerate}[label=\Alph*., leftmargin=*]
\item Il a procédé à une levée de fonds en émettant des actions et des obligations, récoltant plus de 3 200 bitcoins
\item Il a obtenu un prêt bancaire traditionnel de 16 000 dollars auprès d'une banque roumaine
\item Il a reçu un investissement de Roger Ver et d'autres business angels de l'écosystème Bitcoin
\item Il a vendu ses options d'achat et de vente de son service MPOE pour financer le développement
\end{enumerate}

\bigskip


\subsection*{Question 49}
Quel était le principal avantage de la vérification de paiement simplifiée (SPV) implémentée dans BitCoinJ et utilisée par les portefeuilles légers comme le Bitcoin Wallet for Android?

\begin{enumerate}[label=\Alph*., leftmargin=*]
\item Elle chiffrait automatiquement les clés privées avec un algorithme AES-256
\item Elle mélangeait automatiquement les bitcoins pour améliorer la confidentialité des transactions
\item Elle générait des adresses personnalisées commençant par des caractères spécifiques
\item Elle permettait de vérifier les transactions sans télécharger l'intégralité de la chaîne de blocs
\end{enumerate}

\bigskip


\subsection*{Question 50}
Quel était le principal problème de sécurité que les portefeuilles déterministes cherchaient à résoudre par rapport au système de génération aléatoire des clés du logiciel principal Bitcoin?

\begin{enumerate}[label=\Alph*., leftmargin=*]
\item La vulnérabilité aux attaques par force brute due à la faible entropie des générateurs de nombres aléatoires
\item L'exposition des clés privées lors de leur transmission sur le réseau Bitcoin pendant les transactions
\item L'impossibilité de chiffrer les clés privées stockées localement sur l'ordinateur de l'utilisateur
\item La nécessité de faire des sauvegardes répétées du fichier wallet.dat à chaque génération de nouvelle clé
\end{enumerate}

\bigskip


\subsection*{Question 51}
Quel était le principal inconvénient du système BCCAPI de Jan Møller qui a motivé Thomas Voegtlin à développer Electrum?

\begin{enumerate}[label=\Alph*., leftmargin=*]
\item Les frais de transaction étaient trop élevés et le temps de synchronisation était excessivement long
\item Le code de l'infrastructure serveur n'était pas public et le système gardait en mémoire les transactions liées aux portefeuilles
\item Le système ne supportait que les téléphones Android et était incompatible avec les ordinateurs de bureau
\item La génération déterministe des clés privées était défaillante et causait des pertes de fonds
\end{enumerate}

\bigskip


\subsection*{Question 52}
Quelle était la principale innovation technique introduite par les portefeuilles déterministes comme celui d'Electrum par rapport aux portefeuilles traditionnels?

\begin{enumerate}[label=\Alph*., leftmargin=*]
\item L'intégration automatique de services de mélange de bitcoins pour améliorer la confidentialité des transactions
\item La possibilité de se connecter directement aux serveurs Bitcoin sans télécharger la blockchain complète
\item La dérivation de toutes les clés privées à partir d'une seule graine (seed) permettant une sauvegarde simplifiée
\item La génération de nouvelles adresses publiques sans avoir besoin de connaître les clés privées correspondantes
\end{enumerate}

\bigskip


\subsection*{Question 53}
Quelle était la principale caractéristique technique qui distinguait le service Instawallet des autres applications dépositaires de l'époque?

\begin{enumerate}[label=\Alph*., leftmargin=*]
\item Il proposait un chiffrement côté client des clés privées avant envoi au serveur
\item L'accès se faisait par une URL unique générée automatiquement sans nécessiter d'inscription préalable
\item Il intégrait nativement la vérification de paiement simplifiée pour plus de sécurité
\item Il utilisait un système de double authentification avec clé privée et mot de passe
\end{enumerate}

\bigskip


\subsection*{Question 54}
Quelle était la principale méthode utilisée par les bitcoins physiques de Casascius pour empêcher l'accès frauduleux à la clé privée?

\begin{enumerate}[label=\Alph*., leftmargin=*]
\item Un chiffrement AES-256 de la clé privée avec un mot de passe unique gravé sur la pièce
\item Un hologramme personnalisé recouvrant la clé privée qui s'endommage lors de toute tentative d'accès
\item Un code QR invisible révélé seulement sous lumière ultraviolette spéciale
\item Une signature cryptographique vérifiable uniquement par le serveur central de Casascius
\end{enumerate}

\bigskip


\subsection*{Question 55}
Quel était le principal objectif de Gregory Maxwell en proposant les portefeuilles déterministes de type 2 par rapport au type 1?

\begin{enumerate}[label=\Alph*., leftmargin=*]
\item Permettre la dérivation des adresses à partir d'une clé publique maîtresse sans exposer la graine
\item Augmenter la sécurité cryptographique en utilisant des fonctions de hachage plus robustes
\item Réduire la taille de stockage nécessaire pour sauvegarder l'ensemble des clés privées
\item Améliorer la vitesse de génération des clés privées en utilisant des algorithmes plus efficaces
\end{enumerate}

\bigskip


\subsection*{Question 56}
Quelle était la principale raison pour laquelle Roman Sterlingov a conçu Bitcoin Fog différemment de BitLaundry en imposant un système de solde local?

\begin{enumerate}[label=\Alph*., leftmargin=*]
\item Pour réduire les coûts opérationnels en évitant de devoir gérer des adresses à usage unique
\item Pour permettre aux utilisateurs de conserver leurs bitcoins de manière permanente sur la plateforme comme un portefeuille
\item Pour éviter les frais de transaction en traitant tous les mélanges en interne sans utiliser la blockchain
\item Pour permettre des retraits différés en plusieurs étapes vers différentes adresses sur une période de 6 à 96 heures
\end{enumerate}

\bigskip


\subsection*{Question 57}
Quel était le principal inconvénient du service StrongCoin développé par DogIsland qui limitait son adoption par rapport à My Wallet de Blockchain.info?

\begin{enumerate}[label=\Alph*., leftmargin=*]
\item StrongCoin ne chiffrait pas les clés privées avant de les envoyer au serveur contrairement à My Wallet
\item StrongCoin n'offrait pas la possibilité de visualiser les détails des transactions contrairement à l'explorateur intégré de Blockchain.info
\item StrongCoin prélevait des frais de 1% sur chaque montant envoyé, allant jusqu'à 1 BTC par transaction
\item StrongCoin nécessitait le téléchargement d'une application mobile alors que My Wallet fonctionnait directement dans le navigateur
\end{enumerate}

\bigskip


\subsection*{Question 58}
Pourquoi Ben Reeves a-t-il développé l'explorateur de blocs Blockchain.info plutôt que d'utiliser les solutions existantes comme BBE ou ABE?

\begin{enumerate}[label=\Alph*., leftmargin=*]
\item Pour développer une alternative open-source gratuite aux explorateurs commerciaux existants
\item Pour créer un explorateur incluant les blocs orphelins et estimant le volume réel des BTC transactés
\item Pour créer le premier explorateur compatible avec les portefeuilles déterministes et la dérivation de clés
\item Pour proposer une interface multilingue supportant l'espéranto et le système de numération tonale
\end{enumerate}

\bigskip


\subsection*{Question 59}
Quelle était la principale innovation technique de l'interface BCCAPI de Jan Møller par rapport aux autres solutions de portefeuilles de l'époque?

\begin{enumerate}[label=\Alph*., leftmargin=*]
\item La possibilité de gérer plusieurs comptes Bitcoin simultanément dans une seule interface utilisateur
\item L'implémentation de la génération déterministe des clés privées éliminant le besoin de sauvegardes régulières
\item L'utilisation de la vérification de paiement simplifiée pour réduire la bande passante nécessaire
\item Le stockage chiffré des clés privées sur des serveurs distants avec authentification par mot de passe
\end{enumerate}

\bigskip


\subsection*{Question 60}
Quelle était la principale motivation de Mike Caldwell (Casascius) pour créer des bitcoins physiques plutôt que de se contenter de solutions numériques existantes?

\begin{enumerate}[label=\Alph*., leftmargin=*]
\item Créer une alternative commerciale rentable aux portefeuilles papier qui étaient trop fragiles pour un usage quotidien
\item Permettre l'échange de bitcoins en personne sans nécessiter d'accès à Internet ou d'appareil électronique
\item Servir d'outil pédagogique pour aider le grand public à visualiser une monnaie virtuelle avec un objet tangible fonctionnel
\item Résoudre les problèmes de sécurité liés au stockage numérique des clés privées sur les ordinateurs connectés
\end{enumerate}

\bigskip


\subsection*{Question 61}
Quel était le principal problème du modèle de rémunération proportionnel simple initialement utilisé par la coopérative de Slush?

\begin{enumerate}[label=\Alph*., leftmargin=*]
\item Les mineurs pouvaient pratiquer le pool hopping en se déconnectant après avoir soumis un certain nombre de parts pour obtenir environ 28% de bitcoins en plus
\item Le calcul des parts était trop complexe et nécessitait une puissance de calcul supplémentaire qui réduisait la rentabilité du minage
\item La variance était trop importante pour les gros mineurs qui préféraient un système de paiement immédiat comme le Pay Per Share
\item Les frais de commission étaient trop élevés à 2% par bloc miné, ce qui décourageait les petits mineurs de participer à la coopérative
\end{enumerate}

\bigskip


\subsection*{Question 62}
Quel protocole de communication a été développé par Marek Palatinus (Slush) en septembre 2012 pour améliorer les échanges entre les mineurs et les coopératives?

\begin{enumerate}[label=\Alph*., leftmargin=*]
\item Le protocole getblocktemplate (GBT)
\item Le protocole long polling
\item Le protocole Stratum
\item Le protocole getmemorypool (GMP)
\end{enumerate}

\bigskip


\subsection*{Question 63}
Quelle méthode de rémunération alternative au modèle proportionnel a été développée par Luke-Jr pour sa coopérative Eligius afin de résoudre les problèmes de variance sans prendre de risque financier excessif?

\begin{enumerate}[label=\Alph*., leftmargin=*]
\item Le modèle de score géométrique double
\item Le modèle Capped Pay Per Share with Recent Backpay (CPPSRB)
\item Le modèle Pay Per Last N Shares (PPLNS)
\item Le modèle Shared Maximum Pay Per Share (SMPPS)
\end{enumerate}

\bigskip


\subsection*{Question 64}
Quelle innovation technique a été introduite par la coopérative Eligius de Luke-Jr le 14 juin 2011 pour améliorer la transparence du fonctionnement des coopératives de minage?

\begin{enumerate}[label=\Alph*., leftmargin=*]
\item La signature de la coinbase consistant à écrire le nom de la coopérative dans le champ d'entrée de la transaction de récompense
\item La création d'un système de sharechain décentralisée permettant aux mineurs de se connecter en pair à pair
\item L'implémentation du modèle de rémunération Pay Per Last N Shares (PPLNS) pour éviter le pool hopping
\item Le développement du protocole de communication Stratum pour réduire la consommation de bande passante
\end{enumerate}

\bigskip


\subsection*{Question 65}
Quel était le pourcentage approximatif de la puissance de calcul du réseau Bitcoin contrôlé par la coopérative Deepbit au pic de sa domination en août 2011?

\begin{enumerate}[label=\Alph*., leftmargin=*]
\item Environ 75% de la puissance totale du réseau
\item Environ 35% de la puissance totale du réseau
\item Environ 28% de la puissance totale du réseau
\item Environ 53% de la puissance totale du réseau
\end{enumerate}

\bigskip


\subsection*{Question 66}
Quel était le principal avantage du système P2Pool développé par Forrest Voight par rapport aux coopératives de minage traditionnelles?

\begin{enumerate}[label=\Alph*., leftmargin=*]
\item Il offrait des frais de commission plus bas que les autres coopératives avec seulement 1% de commission
\item Il permettait un minage décentralisé sans dépendre d'un opérateur central grâce à une sharechain maintenue en pair à pair
\item Il utilisait un algorithme de hachage plus efficace qui permettait d'augmenter la rentabilité de 30%
\item Il permettait de miner plusieurs cryptomonnaies simultanément sans perte de performance
\end{enumerate}

\bigskip


\subsection*{Question 67}
Pourquoi la méthode Pay Per Last N Shares (PPLNS) développée par Wuked pour Mineco.in était-elle considérée comme une amélioration par rapport au modèle proportionnel classique?

\begin{enumerate}[label=\Alph*., leftmargin=*]
\item Elle utilisait une fenêtre de N parts indépendante du cycle de production des blocs pour calculer les récompenses
\item Elle combinait le modèle proportionnel avec un système de tampon collectif pour éviter les pertes
\item Elle appliquait un score décroissant géométriquement pour favoriser les parts les plus récentes
\item Elle permettait de payer immédiatement chaque part fournie selon son degré de difficulté
\end{enumerate}

\bigskip


\subsection*{Question 68}
Quelle était la particularité technique fondamentale du système P2Pool qui le distinguait des autres coopératives de minage?

\begin{enumerate}[label=\Alph*., leftmargin=*]
\item Il permettait le minage simultané de Bitcoin et Namecoin sans perte de performance
\item Il utilisait une sharechain maintenue en pair à pair par les participants sans serveur central
\item Il utilisait le protocole Stratum pour optimiser la communication entre mineurs
\item Il appliquait exclusivement un modèle de rémunération Pay Per Share avec commission réduite
\end{enumerate}

\bigskip


\subsection*{Question 69}
Quelle était la principale motivation économique qui a poussé les mineurs individuels à rejoindre les coopératives de minage à partir de l'automne 2010?

\begin{enumerate}[label=\Alph*., leftmargin=*]
\item L'accès exclusif aux nouvelles technologies de minage GPU réservées aux membres des coopératives
\item La réduction de la variance des revenus car les petits mineurs ne pouvaient plus espérer trouver un bloc dans un délai raisonnable
\item La possibilité d'obtenir des commissions plus élevées en mutualisant les coûts d'électricité entre participants
\item La garantie de recevoir des bitcoins même sans contribuer activement au processus de minage
\end{enumerate}

\bigskip


\subsection*{Question 70}
Quelle était la difficulté initiale de la sharechain de P2Pool par rapport à la difficulté du réseau Bitcoin principal?

\begin{enumerate}[label=\Alph*., leftmargin=*]
\item 300 fois inférieure à la difficulté du réseau principal
\item 1000 fois inférieure à la difficulté du réseau principal
\item 150 fois inférieure à la difficulté du réseau principal
\item 600 fois inférieure à la difficulté du réseau principal
\end{enumerate}

\bigskip


\subsection*{Question 71}
Quelle était la principale raison pour laquelle Luke-Jr a abandonné l'inclusion de messages religieux dans les blocs d'Eligius au début septembre 2011?

\begin{enumerate}[label=\Alph*., leftmargin=*]
\item Des problèmes techniques causés par la longueur excessive des prières latines dans les transactions
\item Une directive officielle des développeurs de Bitcoin interdisant les messages personnels dans la coinbase
\item Les critiques de la communauté concernant l'inclusion de contenu religieux sans le consentement des mineurs participants
\item Une menace de poursuites judiciaires de la part d'organisations religieuses concurrentes
\end{enumerate}

\bigskip


\subsection*{Question 72}
Quel était l'intervalle de temps moyen visé entre les parts dans la sharechain du système P2Pool développé par Forrest Voight?

\begin{enumerate}[label=\Alph*., leftmargin=*]
\item 30 secondes
\item 5 secondes
\item 10 minutes
\item 2 minutes
\end{enumerate}

\bigskip


\subsection*{Question 73}
Quelle était la principale préoccupation de Russell O'Connor concernant la proposition OP_EVAL?

\begin{enumerate}[label=\Alph*., leftmargin=*]
\item L'impossibilité d'analyser statiquement le script sans l'exécuter et les risques de récursion
\item Le manque de compatibilité avec les anciennes versions du logiciel Bitcoin
\item La complexité excessive du processus de signalement par les mineurs
\item L'augmentation trop importante de la taille des blocs et des transactions
\end{enumerate}

\bigskip


\subsection*{Question 74}
Quelle était la principale raison pour laquelle Gavin Andresen a finalement privilégié Pay to Script Hash (P2SH) par rapport à OP_EVAL?

\begin{enumerate}[label=\Alph*., leftmargin=*]
\item P2SH permettait une analyse statique du script sans l'exécuter, contrairement à OP_EVAL qui nécessitait l'exécution
\item P2SH évitait complètement le risque de scission de chaîne contrairement aux autres propositions
\item P2SH bénéficiait du soutien unanime de tous les développeurs principaux de Bitcoin à l'époque
\item P2SH était plus facile à implémenter techniquement et demandait moins de modifications du code existant
\end{enumerate}

\bigskip


\subsection*{Question 75}
Quel était le principal défaut technique d'OP_EVAL qui a conduit à son abandon au profit de Pay to Script Hash?

\begin{enumerate}[label=\Alph*., leftmargin=*]
\item L'augmentation excessive de la taille des adresses Bitcoin rendant les transactions trop coûteuses
\item La limitation du nombre de participants à seulement 3 signataires maximum dans les transactions
\item L'incompatibilité totale avec les anciennes versions du logiciel nécessitant un hard fork
\item L'impossibilité d'analyser statiquement le script sans l'exécuter, empêchant de compter les opérateurs de signature
\end{enumerate}

\bigskip


\subsection*{Question 76}
Pourquoi l'activation de Pay to Script Hash (P2SH) a-t-elle été retardée plusieurs fois au début de 2012?

\begin{enumerate}[label=\Alph*., leftmargin=*]
\item La découverte de bugs critiques dans l'implémentation du code P2SH nécessitant des corrections majeures
\item Les problèmes de compatibilité avec les anciennes versions du logiciel Bitcoin-Qt et bitcoind
\item Le manque de soutien des coopératives de minage, notamment DeepBit qui contrôlait 30% de la puissance de calcul
\item L'opposition unanime des développeurs principaux qui préféraient la solution OP_EVAL originale
\end{enumerate}

\bigskip


\subsection*{Question 77}
Quelle était la principale innovation technique de Pay to Script Hash (P2SH) par rapport aux transactions multisignatures traditionnelles?

\begin{enumerate}[label=\Alph*., leftmargin=*]
\item Il permettait de payer une empreinte de script et de fournir le script complet lors du déverrouillage des fonds
\item Il permettait d'exécuter des boucles et des opérations récursives dans le langage de script Bitcoin
\item Il remplaçait complètement le système OP_CHECKMULTISIG par un nouveau mécanisme de validation
\item Il augmentait le nombre maximum de signataires possibles de 3 à 15 participants dans une transaction
\end{enumerate}

\bigskip


\subsection*{Question 78}
Quel était le principal mécanisme utilisé pour activer Pay to Script Hash (P2SH) sans créer de scission de chaîne?

\begin{enumerate}[label=\Alph*., leftmargin=*]
\item Une modification automatique du protocole après une date prédéfinie
\item Le signalement par les mineurs incluant la chaîne de caractères '/P2SH/' dans les blocs produits
\item Un vote direct de la communauté organisé sur le forum BitcoinTalk
\item L'approbation unanime de tous les développeurs principaux de Bitcoin
\end{enumerate}

\bigskip


\subsection*{Question 79}
Quel était le rôle du BIP 30 dans le développement de Bitcoin en mars 2012?

\begin{enumerate}[label=\Alph*., leftmargin=*]
\item Il définissait le nouveau format d'adresses courtes pour les transactions Pay to Script Hash
\item Il établissait le mécanisme de signalement des mineurs pour l'activation des soft forks
\item Il introduisait l'opérateur OP_CHECKMULTISIG pour les signatures multipartites dans le protocole Bitcoin
\item Il interdisait aux blocs futurs de contenir des transactions avec des identifiants identiques à des transactions antérieures non entièrement dépensées
\end{enumerate}

\bigskip


\subsection*{Question 80}
Quelle était la principale raison technique pour laquelle Gavin Andresen a proposé d'utiliser les opérateurs OP_NOP pour implémenter les changements de protocole?

\begin{enumerate}[label=\Alph*., leftmargin=*]
\item Pour permettre une modification rétrocompatible évitant une scission de chaîne
\item Pour réduire la taille des scripts et optimiser l'espace de stockage
\item Pour accélérer le processus de validation des transactions multisignatures
\item Pour corriger automatiquement le bug présent dans OP_CHECKMULTISIG
\end{enumerate}

\bigskip


\subsection*{Question 81}
Quelle était la principale différence entre le système de développement de Bitcoin après le départ de Satoshi et celui mis en place par Satoshi lui-même?

\begin{enumerate}[label=\Alph*., leftmargin=*]
\item Le développement est passé d'un système ouvert à un système fermé contrôlé par une seule personne
\item Le développement s'est centralisé autour d'une entreprise privée dirigée par Gavin Andresen
\item Le développement est devenu collaboratif avec plusieurs mainteneurs ayant accès en écriture au dépôt GitHub
\item Le développement a adopté un modèle de vote démocratique où chaque utilisateur avait une voix
\end{enumerate}

\bigskip


\subsection*{Question 82}
Quelle était la principale raison pour laquelle Tycho de Deepbit s'opposait initialement au déploiement de Pay to Script Hash?

\begin{enumerate}[label=\Alph*., leftmargin=*]
\item Il considérait qu'il n'y avait pas de consensus entre les développeurs et que la solution était trop bricolée
\item Il craignait que P2SH ne réduise les frais de transaction et donc les revenus de sa coopérative
\item Il voulait imposer sa propre proposition technique alternative basée sur OP_CHECKHASHVERIFY
\item Il préférait maintenir sa position dominante sur le réseau avec 30% de la puissance de calcul
\end{enumerate}

\bigskip


\subsection*{Question 83}
Quelle était la principale caractéristique technique qui distinguait le framework Qt utilisé par Wladimir van der Laan de la solution graphique précédente de Bitcoin?

\begin{enumerate}[label=\Alph*., leftmargin=*]
\item Il utilisait un système de validation des blocs plus rapide que l'implémentation précédente
\item Il permettait l'exécution de scripts multisignatures directement dans l'interface graphique
\item Il intégrait nativement le support des transactions Pay to Script Hash dans l'interface
\item Il remplaçait la bibliothèque wxWindows utilisée par Satoshi par le framework graphique Qt
\end{enumerate}

\bigskip


\subsection*{Question 84}
Quel était le principal problème de sécurité découvert par Russell O'Connor en février 2012 qui a nécessité la création du BIP 30?

\begin{enumerate}[label=\Alph*., leftmargin=*]
\item Un bug de récursion dans l'opérateur OP_CHECKMULTISIG qui permettait des boucles infinies dans les scripts
\item La possibilité de supprimer toutes les transactions dérivées d'une transaction de récompense en double via une réorganisation de chaîne
\item Un défaut dans le système de signalement des mineurs qui permettait de falsifier le support pour les BIP
\item Une vulnérabilité dans l'interface Bitcoin-Qt qui exposait les clés privées lors des transactions multisignatures
\end{enumerate}

\bigskip


\subsection*{Question 85}
Quel procédé technique Namecoin a-t-il introduit pour permettre aux mineurs de Bitcoin de participer à son réseau sans effort supplémentaire?

\begin{enumerate}[label=\Alph*., leftmargin=*]
\item L'algorithme Scrypt qui rend le minage résistant aux processeurs graphiques
\item Le système de nœuds de confiance qui valident les transactions en parallèle
\item La preuve d'enjeu hybride qui combine possession d'unités et puissance de calcul
\item Le minage combiné (merged mining) qui réutilise le travail fourni sur la chaîne Bitcoin
\end{enumerate}

\bigskip


\subsection*{Question 86}
Quelle était la principale innovation technique introduite par Tenebrix pour modifier l'équilibre du minage par rapport à Bitcoin?

\begin{enumerate}[label=\Alph*., leftmargin=*]
\item L'implémentation d'un système de preuve d'enjeu hybride pour réduire la consommation énergétique
\item La création d'un temps de bloc de 15 secondes pour accélérer drastiquement les confirmations
\item L'introduction de nœuds de confiance pour protéger le réseau contre les attaques des 51%
\item L'utilisation de la fonction Scrypt pour rendre le minage résistant aux GPU et favorable aux CPU
\end{enumerate}

\bigskip


\subsection*{Question 87}
Quelle était la principale différence entre Ixcoin et I0coin lors de leur lancement en août 2011?

\begin{enumerate}[label=\Alph*., leftmargin=*]
\item Ixcoin était compatible avec le minage combiné tandis qu'I0coin nécessitait un minage séparé
\item Ixcoin avait un temps de bloc de 10 minutes tandis qu'I0coin avait un temps de bloc de 2,5 minutes
\item Ixcoin avait été préminiée avec plus de 580 000 unités tandis qu'I0coin n'avait aucun bloc prégénéré
\item Ixcoin utilisait l'algorithme Scrypt pour le minage tandis qu'I0coin utilisait SHA-256
\end{enumerate}

\bigskip


\subsection*{Question 88}
Quel problème technique majeur a affecté le lancement de Litecoin le 13 octobre 2011?

\begin{enumerate}[label=\Alph*., leftmargin=*]
\item Plus de 10 000 blocs ont été minés en moins de 24 heures à cause d'une difficulté initiale trop basse
\item Charlie Lee avait prépminé secrètement 500 000 litecoins avant le lancement public
\item L'algorithme Scrypt ne fonctionnait pas correctement et générait des blocs invalides
\item Le réseau a subi une attaque des 51% qui a compromis la sécurité de la blockchain
\end{enumerate}

\bigskip


\subsection*{Question 89}
Selon Sunny King, quelle était la principale avancée révolutionnaire introduite par Bitcoin qui a permis l'émergence des cryptomonnaies?

\begin{enumerate}[label=\Alph*., leftmargin=*]
\item L'implémentation d'un réseau pair-à-pair
\item L'utilisation de la fonction de hachage SHA-256
\item L'introduction d'un registre public décentralisé
\item La création d'un système de preuve de travail
\end{enumerate}

\bigskip


\subsection*{Question 90}
Quelle expression péjorative est apparue dans la communauté Bitcoin pour critiquer la prolifération excessive de nouvelles cryptomonnaies alternatives?

\begin{enumerate}[label=\Alph*., leftmargin=*]
\item Yet another cryptocurrency
\item Altcoin multiplication disorder
\item Bitcoin copycat syndrome
\item Cryptocurrency spam attack
\end{enumerate}

\bigskip


\subsection*{Question 91}
Quel était le principal objectif du triangle de Zooko théorisé en 2001 par Bryce Wilcox dans le contexte des systèmes de noms?

\begin{enumerate}[label=\Alph*., leftmargin=*]
\item Montrer que trois types d'autorités centrales doivent valider simultanément chaque nom de domaine enregistré
\item Prouver que trois serveurs DNS suffisent pour assurer la redondance et la sécurité d'un système de noms
\item Établir que trois algorithmes cryptographiques sont nécessaires pour sécuriser efficacement les noms de domaine
\item Démontrer qu'un identifiant ne peut réunir que deux des trois propriétés, sécurité, décentralisation et intelligibilité humaine
\end{enumerate}

\bigskip


\subsection*{Question 92}
Quelle était la particularité du modèle économique de Groupcoin/Devcoin lancé en juillet 2011 par rapport aux autres cryptomonnaies alternatives de l'époque?

\begin{enumerate}[label=\Alph*., leftmargin=*]
\item 90% des nouvelles unités créées étaient distribuées aux développeurs de logiciels libres, seuls 10% allant aux mineurs
\item Le système fonctionnait uniquement par preuve d'enjeu sans aucune récompense pour les mineurs
\item 100% des unités étaient préminées et vendues directement sur les plateformes d'échange sans minage
\item Les unités étaient distribuées équitablement entre tous les utilisateurs inscrits sur le forum BitcoinTalk
\end{enumerate}

\bigskip


\subsection*{Question 93}
Quelle action controversée Luke-Jr a-t-il entreprise contre Coiledcoin le 6 janvier 2012 pour protéger l'écosystème Bitcoin?

\begin{enumerate}[label=\Alph*., leftmargin=*]
\item Il a lancé une campagne de spam sur les forums pour discréditer techniquement le projet Coiledcoin
\item Il a modifié la page Wikipedia de Coiledcoin pour la qualifier publiquement d'arnaque pyramidale
\item Il a réalisé une attaque de censure fatale contre Coiledcoin à l'aide de sa coopérative de minage Eligius
\item Il a déposé un avis de retrait DMCA auprès de GitHub pour violation de droits d'auteur sur le code Bitcoin
\end{enumerate}

\bigskip


\subsection*{Question 94}
Quelle était la principale motivation derrière la création du système BitDNS proposé par Appamatto en novembre 2010?

\begin{enumerate}[label=\Alph*., leftmargin=*]
\item Lancer une plateforme de change décentralisée pour faciliter les échanges entre différentes cryptomonnaies
\item Développer une cryptomonnaie plus rapide que Bitcoin avec des temps de bloc réduits pour les transactions
\item Implémenter un système de preuve d'enjeu pour réduire la consommation énergétique du minage de Bitcoin
\item Créer un système de noms de domaine décentralisé résistant à la censure pour remplacer le contrôle centralisé de l'ICANN
\end{enumerate}

\bigskip


\subsection*{Question 95}
Quel était le principal défaut technique qui a causé des problèmes similaires lors du lancement de Tenebrix et de Litecoin?

\begin{enumerate}[label=\Alph*., leftmargin=*]
\item Le temps de bloc était trop rapide ce qui créait automatiquement des blocs orphelins en cascade
\item Le système de récompenses était défaillant et distribuait plus d'unités que prévu par bloc miné
\item La fonction de hachage Scrypt était mal implémentée et causait des erreurs de validation des blocs
\item L'algorithme d'ajustement de difficulté était limité et ne pouvait multiplier ou diviser la difficulté que par 4 maximum
\end{enumerate}

\bigskip


\subsection*{Question 96}
Quelle était la principale caractéristique technique de la fonction de hachage Scrypt implémentée par ArtForz dans Tenebrix pour modifier l'équilibre du minage?

\begin{enumerate}[label=\Alph*., leftmargin=*]
\item Une forte demande de mémoire qui réduisait drastiquement l'efficacité des processeurs graphiques
\item Une réduction du temps de calcul qui permettait des blocs plus rapides
\item Un système de hachage parallèle qui optimisait l'utilisation des processeurs multicœurs
\item Un algorithme de double hachage qui augmentait la sécurité cryptographique du réseau
\end{enumerate}

\bigskip


\subsection*{Question 97}
Quelle était la principale justification économique de Ross Ulbricht pour modifier le système de commissions de Silk Road en janvier 2012?

\begin{enumerate}[label=\Alph*., leftmargin=*]
\item Réduire les frais à un taux fixe de 3,7% pour concurrencer directement les plateformes légales comme eBay et Amazon
\item Augmenter uniformément les commissions à 15% pour tous les vendeurs afin de financer l'expansion internationale de la plateforme
\item Instaurer des commissions variables selon le type de produit vendu pour décourager certaines catégories de drogues dangereuses
\item Remplacer un taux fixe de 6,23% par un système dégressif favorisant les grosses transactions plutôt que les petites ventes
\end{enumerate}

\bigskip


\subsection*{Question 98}
Quelle était la principale innovation technique introduite par LocalBitcoins pour sécuriser les échanges de bitcoins contre espèces entre particuliers?

\begin{enumerate}[label=\Alph*., leftmargin=*]
\item Un système de géolocalisation automatique permettant de vérifier la présence physique des deux parties lors de l'échange
\item Un système de dépôt fiduciaire avec déblocage par SMS permettant de conserver les bitcoins jusqu'à confirmation de réception des espèces
\item Un protocole de chiffrement avancé intégré au navigateur Tor pour protéger l'identité des utilisateurs pendant les transactions
\item Un mécanisme de validation par empreinte digitale obligatoire pour tous les participants aux échanges en personne
\end{enumerate}

\bigskip


\subsection*{Question 99}
Quelle était la principale stratégie économique de BitPay qui la différenciait fondamentalement de BitInstant dans le traitement des bitcoins?

\begin{enumerate}[label=\Alph*., leftmargin=*]
\item BitPay facilitait les échanges de bitcoins entre particuliers via un système de dépôt fiduciaire par SMS
\item BitPay permettait aux clients de déposer des espèces dans des banques physiques pour obtenir des bitcoins instantanément
\item BitPay achetait les bitcoins des clients contre des espèces puis les revendait sur les places de marché clandestines
\item BitPay avançait les dollars aux commerçants et se chargeait de vendre les bitcoins plus tard pour éviter la volatilité
\end{enumerate}

\bigskip


\subsection*{Question 100}
Quelle référence culturelle Ross Ulbricht a-t-il choisie pour son nouveau pseudonyme en février 2012, et quelle était la signification stratégique de ce choix?

\begin{enumerate}[label=\Alph*., leftmargin=*]
\item Le roman Treasure Island où le capitaine Roberts symbolise l'autorité maritime incontestée sur les océans
\item Le film The Princess Bride où Dread Pirate Roberts est un titre transmissible permettant de masquer l'identité réelle du pirate
\item Le livre Black Sails où Roberts incarne la résistance contre les monopoles commerciaux des grandes puissances
\item La série Pirates of the Caribbean où Dread Pirate Roberts représente la liberté absolue face aux autorités coloniales
\end{enumerate}

\bigskip


\subsection*{Question 101}
Quel pourcentage approximatif de l'activité économique totale de la blockchain Bitcoin représentait Silk Road durant l'année 2012?

\begin{enumerate}[label=\Alph*., leftmargin=*]
\item Moins de 5 % car dominé par Mt. Gox et les échanges
\item Plus de 50 % jusqu'à l'émergence de BitPay en fin d'année
\item Entre 30 et 40 % de manière constante toute l'année
\item Entre 10 et 20 % avec des pics en février et septembre
\end{enumerate}

\bigskip


\subsection*{Question 102}
Quelle était la principale innovation du schéma d'URI bitcoin standardisé par le BIP 21 pour faciliter les paiements commerciaux?

\begin{enumerate}[label=\Alph*., leftmargin=*]
\item Il permettait de valider instantanément les transactions sans attendre les confirmations du réseau
\item Il permettait de crypter automatiquement les transactions pour protéger l'identité des commerçants
\item Il permettait d'encoder dans un code QR l'adresse du destinataire, le montant à payer et la description de la transaction
\item Il permettait de convertir automatiquement les bitcoins en dollars lors du scan du code QR
\end{enumerate}

\bigskip


\subsection*{Question 103}
Selon l'étude de Nicolas Christin publiée en août 2012, quelle était l'évolution du volume de vente journalier de Silk Road entre mars et juillet 2012?

\begin{enumerate}[label=\Alph*., leftmargin=*]
\item Le volume a doublé de 8 000 BTC en mars à 15 000 BTC en mai, puis a reculé à 11 000 BTC en juillet à cause de la hausse du prix du bitcoin
\item Le volume a chuté de 12 000 BTC en mars à 8 000 BTC en juillet suite aux premières enquêtes policières sur la plateforme
\item Le volume est resté stable autour de 10 000 BTC durant toute la période en raison de la stagnation du prix du bitcoin à 5 dollars
\item Le volume a triplé de façon constante de 5 000 BTC en mars à 15 000 BTC en juillet grâce à l'expansion internationale de la plateforme
\end{enumerate}

\bigskip


\subsection*{Question 104}
Quelle était la principale raison économique qui a motivé Robert Faiella (BTCKing) à utiliser spécifiquement le service de dépôt d'espèces de BitInstant pour son activité de revente de bitcoins sur Silk Road?

\begin{enumerate}[label=\Alph*., leftmargin=*]
\item Il était le seul service capable de traiter des volumes supérieurs à 50 000 dollars par transaction
\item Il offrait les frais de commission les plus bas du marché avec seulement 2% de frais totaux
\item Il permettait d'acheter des bitcoins de manière totalement confidentielle sans laisser de trace de virement bancaire
\item Il proposait une conversion automatique en monnaies étrangères pour les clients internationaux
\end{enumerate}

\bigskip


\subsection*{Question 105}
Quelle était la principale justification de Charlie Shrem pour continuer à faciliter les transactions de BTCKing malgré les réticences de son cofondateur Gareth Nelson?

\begin{enumerate}[label=\Alph*., leftmargin=*]
\item Les transactions de BTCKing représentaient moins de 5% du volume total de BitInstant et ne posaient donc pas de risque significatif
\item BTCKing avait obtenu une licence officielle de change d'argent en Floride et opérait dans un cadre réglementaire strict
\item BTCKing n'avait enfreint aucune loi et Silk Road elle-même n'était pas illégale, tout en générant de gros bénéfices pour BitInstant
\item BTCKing respectait scrupuleusement les limites de dépôt de 1000$ et utilisait des adresses email distinctes pour chaque transaction
\end{enumerate}

\bigskip


\subsection*{Question 106}
Quelle était la principale conséquence de l'étude de Nicolas Christin publiée en août 2012 sur l'évolution de Silk Road?

\begin{enumerate}[label=\Alph*., leftmargin=*]
\item Elle a provoqué une croissance explosive de la fréquentation de la plateforme en attirant l'attention sur son succès économique
\item Elle a conduit à la fermeture temporaire de la plateforme pour des raisons de sécurité et de discrétion
\item Elle a entraîné une migration massive des utilisateurs vers des plateformes concurrentes moins exposées médiatiquement
\item Elle a forcé Ross Ulbricht à modifier drastiquement le système de commissions pour réduire la visibilité économique
\end{enumerate}

\bigskip


\subsection*{Question 107}
Quel était le modèle économique distinctif de BitPay par rapport aux autres services de l'époque pour gérer la volatilité du bitcoin?

\begin{enumerate}[label=\Alph*., leftmargin=*]
\item BitPay avançait les dollars aux commerçants contre une commission et se chargeait de vendre les bitcoins plus tard
\item BitPay conservait les bitcoins en dépôt fiduciaire et les libérait uniquement après confirmation de la transaction
\item BitPay facturait des frais fixes de 6,23% sur toutes les transactions sans conversion automatique en dollars
\item BitPay permettait aux commerçants d'acheter des bitcoins avec des espèces déposées dans des emplacements bancaires
\end{enumerate}

\bigskip


\subsection*{Question 108}
Quelle était la principale innovation de sécurité introduite par LocalBitcoins durant l'été 2012 pour protéger les échanges entre particuliers?

\begin{enumerate}[label=\Alph*., leftmargin=*]
\item Une assurance automatique couvrant les pertes jusqu'à 1000 dollars par transaction
\item Un système de dépôt fiduciaire avec déblocage par SMS une fois la réception des espèces confirmée
\item Un système de géolocalisation en temps réel pour surveiller les rencontres physiques
\item Une vérification d'identité obligatoire liée aux comptes bancaires des utilisateurs
\end{enumerate}

\bigskip


\subsection*{Question 109}
Quelle innovation technique de SatoshiDICE permettait aux joueurs de vérifier que le casino ne trichait pas?

\begin{enumerate}[label=\Alph*., leftmargin=*]
\item La transparence totale de la blockchain qui rendait tous les calculs de probabilité publiquement visibles
\item Un système de validation décentralisée où chaque transaction était vérifiée par plusieurs nœuds indépendants
\item L'utilisation d'adresses Bitcoin fixes qui garantissaient automatiquement l'équité des tirages aléatoires
\item La publication des empreintes des secrets sur la blockchain et leur dévoilement ultérieur pour permettre le recalcul des résultats
\end{enumerate}

\bigskip


\subsection*{Question 110}
Pourquoi le 15 avril 2011 est-il surnommé le « vendredi noir » par les joueurs de poker en ligne ?

\begin{enumerate}[label=\Alph*., leftmargin=*]
\item Bitcoin a subi sa première chute majeure de prix, affectant tous les sites de jeu en ligne
\item Le FBI a fermé les trois principales plateformes de poker américaines (PokerStars, Full Tilt Poker et Absolute Poker)
\item Une cyberattaque massive a compromis simultanément tous les portefeuilles Bitcoin des joueurs de poker
\item Le gouvernement américain a déclaré Bitcoin illégal pour toutes les transactions de jeu d'argent
\end{enumerate}

\bigskip


\subsection*{Question 111}
Quelle était la principale conséquence du fonctionnement de SatoshiDICE sur l'infrastructure Bitcoin en 2012?

\begin{enumerate}[label=\Alph*., leftmargin=*]
\item Chaque pari générait deux transactions, provoquant une explosion du nombre de transactions quotidiennes sur la blockchain
\item Les calculs de vérification d'équité ralentissaient considérablement la validation des blocs par les mineurs
\item Le système provoquait des congestions en stockant directement les secrets de chiffrement dans chaque bloc
\item L'utilisation massive d'adresses fixes créait des conflits dans le système de routage des transactions
\end{enumerate}

\bigskip


\subsection*{Question 112}
Selon le chapitre, quel était le montant initial payé par Erik Voorhees pour acquérir le concept de SatoshiDICE de Fireduck?

\begin{enumerate}[label=\Alph*., leftmargin=*]
\item 146 bitcoins (environ 730 $ à l'époque)
\item 45 bitcoins (environ 225 $ à l'époque)
\item 40 bitcoins (environ 200 $ à l'époque)
\item 34 500 bitcoins (environ 371 910 $ à l'époque)
\end{enumerate}

\bigskip


\subsection*{Question 113}
Quel était le principal argument technique avancé pour justifier l'utilisation de Bitcoin dans le secteur du jeu d'argent en ligne après 2011?

\begin{enumerate}[label=\Alph*., leftmargin=*]
\item Bitcoin évitait les fluctuations de change qui affectaient les plateformes de jeu internationales
\item Bitcoin réduisait les risques de rétrofacturation et limitait la capacité des États à bloquer les paiements
\item Bitcoin permettait des transactions instantanées sans frais contrairement aux systèmes bancaires traditionnels
\item Bitcoin offrait une anonymat total qui rendait impossible le traçage des transactions de jeu
\end{enumerate}

\bigskip


\subsection*{Question 114}
Quelle était la méthode utilisée par PokerStars, Full Tilt Poker et Absolute Poker pour contourner l'Unlawful Internet Gambling Enforcement Act avant leur fermeture en 2011?

\begin{enumerate}[label=\Alph*., leftmargin=*]
\item Chiffrer toutes les transactions pour les rendre indétectables par les autorités
\item Se faire passer pour des commerçants vendant des bijoux et des balles de golf
\item Créer des filiales bancaires dans des pays sans réglementation sur le jeu
\item Utiliser des serveurs situés dans des paradis fiscaux offshore
\end{enumerate}

\bigskip


\subsection*{Question 115}
Quelle était la proportion approximative des transactions Bitcoin représentée par SatoshiDICE au fil des mois suivant son lancement?

\begin{enumerate}[label=\Alph*., leftmargin=*]
\item Près de 90% des transactions liées au jeu d'argent
\item Approximativement 75% du volume économique total
\item Près de la moitié de toutes les transactions sur la chaîne
\item Environ un quart des transactions quotidiennes
\end{enumerate}

\bigskip


\subsection*{Question 116}
Quel était le mécanisme technique utilisé par SatoshiDICE pour calculer le « nombre chanceux » déterminant l'issue de chaque pari?

\begin{enumerate}[label=\Alph*., leftmargin=*]
\item Génération d'un nombre aléatoire basé sur l'horodatage de la transaction et la clé publique du joueur
\item Sélection des deux premiers octets du condensé HMAC-SHA-512 d'une information secrète et de l'identifiant de transaction
\item Utilisation d'un oracle externe fournissant des nombres pseudo-aléatoires vérifiables par signature cryptographique
\item Calcul d'un hash SHA-256 de l'adresse Bitcoin du joueur combinée avec un sel quotidien prédéfini
\end{enumerate}

\bigskip


\subsection*{Question 117}
Quelle était la principale motivation de Joseph Gleason (Fireduck) pour vendre son casino de bitcoins seulement 10 jours après son lancement en avril 2012?

\begin{enumerate}[label=\Alph*., leftmargin=*]
\item Il manquait de capitaux pour faire face à l'augmentation rapide des mises des joueurs
\item Il appréhendait les conséquences légales et cherchait quelqu'un pour assumer les risques juridiques
\item Il était dépassé par les aspects techniques de la vérification cryptographique des tirages
\item Il voulait se concentrer sur le développement technique plutôt que sur la gestion commerciale
\end{enumerate}

\bigskip


\subsection*{Question 118}
Quelle était la raison principale pour laquelle Erik Voorhees a procédé à une IPO informelle de SatoshiDICE sur MPEx en août 2012?

\begin{enumerate}[label=\Alph*., leftmargin=*]
\item Pour recueillir des fonds d'investisseurs et financer le développement du site grâce à la vente de parts donnant droit à des dividendes
\item Pour créer une cryptomonnaie alternative spécialement dédiée aux jeux d'argent en ligne et concurrencer Bitcoin
\item Pour contourner les réglementations américaines en transférant la propriété légale du site à des investisseurs étrangers
\item Pour réduire sa responsabilité personnelle face aux critiques techniques de la communauté Bitcoin concernant le spam de transactions
\end{enumerate}

\bigskip


\subsection*{Question 119}
Quelle était la particularité technique du système de SatoshiDICE qui le distinguait des autres casinos en ligne de l'époque?

\begin{enumerate}[label=\Alph*., leftmargin=*]
\item L'utilisation d'adresses Bitcoin multisignatures pour sécuriser automatiquement les fonds des joueurs
\item Le caractère « assurément équitable » permettant aux utilisateurs de vérifier mathématiquement l'honnêteté des tirages
\item Une technologie de chiffrement quantique rendant les transactions de jeu totalement anonymes et intraçables
\item Un système de contrats intelligents intégrés permettant l'exécution automatique des paris sans intervention humaine
\end{enumerate}

\bigskip


\subsection*{Question 120}
Quelle était la principale innovation conceptuelle introduite par SatoshiDICE qui a révolutionné le secteur des casinos en ligne utilisant Bitcoin ?

\begin{enumerate}[label=\Alph*., leftmargin=*]
\item Un mécanisme de paris instantanés sans confirmation de transaction nécessaire sur la blockchain Bitcoin
\item Une plateforme décentralisée fonctionnant sans serveur central grâce à des contrats intelligents sur Bitcoin
\item Un système de récompenses automatiques basé sur l'adresse Bitcoin personnelle du joueur sans intervention humaine
\item Un système de jeu « assurément équitable » où les utilisateurs peuvent vérifier mathématiquement que le casino n'a pas triché
\end{enumerate}

\bigskip


\subsection*{Question 121}
Quelle était la particularité du système de taux d'intérêt proposé par Bitcoin Savings & Trust de Trendon Shavers?

\begin{enumerate}[label=\Alph*., leftmargin=*]
\item Un taux fixe de 1% par jour pour tous les investisseurs, soit une multiplication par 37 de la mise en un an
\item Les taux étaient progressifs selon le montant investi, allant de 4,2% par semaine pour 100 BTC à 7% par semaine pour 1000 BTC et plus
\item Des taux d'intérêt variables basés sur les fluctuations du prix du bitcoin, garantissant un rendement en dollars américains
\item Un système de taux dégressifs où les petits investisseurs recevaient 7% par semaine et les gros investisseurs seulement 4,2%
\end{enumerate}

\bigskip


\subsection*{Question 122}
Quelle était la méthode utilisée par le pirate lors du troisième piratage de Bitcoinica en juillet 2012?

\begin{enumerate}[label=\Alph*., leftmargin=*]
\item Il a exploité une clé API contenue dans le code source publié par Amir Taaki qui donnait accès au coffre-fort LastPass
\item Il a utilisé des informations d'identification volées à Zhou Tong pour se connecter aux comptes Mt. Gox de Bitcoinica
\item Il a piraté directement les serveurs de Linode pour accéder aux portefeuilles chauds stockés sur la plateforme d'hébergement
\item Il a compromis le serveur de courrier électronique hébergé par Rackspace pour réinitialiser les mots de passe administrateurs
\end{enumerate}

\bigskip


\subsection*{Question 123}
Quelle était la conséquence directe de l'absence de sauvegarde de la base de données lors du deuxième piratage de Bitcoinica en mai 2012?

\begin{enumerate}[label=\Alph*., leftmargin=*]
\item L'impossibilité totale de rembourser les clients et la fermeture immédiate de la plateforme
\item Le transfert automatique de tous les fonds restants vers les comptes Mt. Gox des utilisateurs
\item La création d'un système de vote communautaire pour déterminer les montants dus à chaque client
\item La mise en place d'un système de demande de remboursement partiel nécessitant la fourniture de preuves par les créanciers
\end{enumerate}

\bigskip


\subsection*{Question 124}
Quelle était la principale vulnérabilité exploitée lors du piratage du service d'hébergement Linode le 1er mars 2012?

\begin{enumerate}[label=\Alph*., leftmargin=*]
\item Un logiciel malveillant installé sur les ordinateurs personnels des clients via des emails de phishing
\item L'accès au portail du service client de l'entreprise permettant de prendre le contrôle de comptes spécifiques
\item Une faille dans le système de chiffrement des portefeuilles Bitcoin stockés sur les serveurs
\item L'utilisation d'une clé API compromise trouvée dans le code source publié accidentellement
\end{enumerate}

\bigskip


\subsection*{Question 125}
Quelle était la principale raison invoquée par Zhou Tong pour justifier sa décision de quitter Bitcoin en mai 2012?

\begin{enumerate}[label=\Alph*., leftmargin=*]
\item Il souhaitait se concentrer sur ses études en Australie et abandonner l'entrepreneuriat
\item Il estimait que la spéculation était un jeu à somme nulle qui ne créait pas de valeur pour la société
\item Il craignait que les autorités réglementaires s'intéressent de trop près à ses activités
\item Il voulait éviter les poursuites judiciaires liées aux piratages successifs de Bitcoinica
\end{enumerate}

\bigskip


\subsection*{Question 126}
Quelle était la caractéristique distinctive de l'escroquerie de Tony76 sur Silk Road le 20 avril 2012 par rapport aux autres fraudes de cette période?

\begin{enumerate}[label=\Alph*., leftmargin=*]
\item Il utilisait une vulnérabilité technique dans le système de paiement Bitcoin de la plateforme
\item Il exploitait sa réputation irréprochable de vendeur établi pour exiger une finalisation anticipée des transactions
\item Il manipulait les cours du bitcoin pour créer des profits artificiels lors des transactions
\item Il profitait d'un piratage informatique des serveurs de Silk Road pour voler les fonds des utilisateurs
\end{enumerate}

\bigskip


\subsection*{Question 127}
Quelle était la conséquence immédiate sur le marché du bitcoin lors de la fermeture de Bitcoin Savings & Trust le 17 août 2012?

\begin{enumerate}[label=\Alph*., leftmargin=*]
\item Le prix a fluctué légèrement à la baisse de 5% puis s'est rapidement rétabli au niveau initial
\item Le prix a augmenté de 50% en raison de la rareté soudaine des bitcoins disponibles sur le marché
\item Le prix est resté stable car les investisseurs avaient anticipé la fermeture du fonds pyramidal
\item Le prix a chuté de moitié en 48 heures, passant de 15,4$ à moins de 8$ avant de se stabiliser autour de 10$
\end{enumerate}

\bigskip


\subsection*{Question 128}
Quelle était la principale innovation de sécurité manquante qui a permis au pirate de voler 24 000 BTC lors du piratage de Bitfloor en septembre 2012?

\begin{enumerate}[label=\Alph*., leftmargin=*]
\item L'absence de chiffrement des clés privées dans les sauvegardes
\item La connexion non sécurisée entre les serveurs de la plateforme
\item L'absence de système d'authentification à deux facteurs
\item L'utilisation d'un mot de passe administrateur trop simple
\end{enumerate}

\bigskip


\subsection*{Question 129}
Quelle était la principale source de revenus revendiquée par Trendon Shavers pour justifier les rendements élevés de Bitcoin Savings & Trust?

\begin{enumerate}[label=\Alph*., leftmargin=*]
\item La revente locale de bitcoins à un taux supérieur à celui de Mt. Gox
\item Le minage de bitcoins via des fermes spécialisées et la plateforme GPUMAX
\item L'arbitrage entre différentes plateformes de change comme Mt. Gox et BTC-e
\item Le prêt de bitcoins à des personnes en contact avec Trendon Shavers
\end{enumerate}

\bigskip


\subsection*{Question 130}
Quelle était la proportion approximative de bitcoins en circulation que représentaient les fonds gérés par Bitcoin Savings & Trust à son apogée?

\begin{enumerate}[label=\Alph*., leftmargin=*]
\item Environ 15% des bitcoins en circulation (plus de 1,5 million BTC sur 9,7 millions)
\item Environ 10% des bitcoins en circulation (environ 1 million BTC sur 9,7 millions)
\item Environ 2% des bitcoins en circulation (environ 200 000 BTC sur 9,7 millions)
\item Environ 5% des bitcoins en circulation (plus de 500 000 BTC sur 9,7 millions)
\end{enumerate}

\bigskip


\subsection*{Question 131}
Quel événement spécifique a directement causé la vulnérabilité exploitée lors du troisième piratage de Bitcoinica en juillet 2012?

\begin{enumerate}[label=\Alph*., leftmargin=*]
\item La suppression accidentelle de la base de données des utilisateurs lors du deuxième piratage
\item La publication publique du code source complet de Bitcoinica par Amir Taaki le 7 juillet
\item La compromission du serveur de courrier électronique hébergé par Rackspace en mai
\item L'accès non autorisé au portail client de Linode par un pirate en mars
\end{enumerate}

\bigskip


\subsection*{Question 132}
Quelle était la principale conséquence structurelle de la multiplication des vols et escroqueries Bitcoin entre 2011 et 2012 sur l'évolution de l'écosystème?

\begin{enumerate}[label=\Alph*., leftmargin=*]
\item L'adoption généralisée du stockage à froid par toutes les plateformes de change pour sécuriser les fonds
\item La création de la Fondation Bitcoin le 27 septembre 2012 pour redorer le blason de Bitcoin et courtiser les autorités
\item L'établissement d'un protocole de sécurité standardisé imposé à tous les acteurs de l'écosystème
\item La mise en place d'un système d'assurance obligatoire pour tous les services Bitcoin par la communauté
\end{enumerate}

\bigskip


\newpage
\section*{Answer Key}

\begin{enumerate}
\item \textbf{Question 1:} D
\\ \textit{Ross Ulbricht utilisait le pseudonyme « Altoid » pour promouvoir Silk Road sur différents forums. Il s'est inscrit sous ce nom sur The Shroomery le 27 janvier 2011, puis sur le forum Bitcoin le 29 janvier 2011. C'est également avec ce pseudonyme qu'il a commis l'erreur fatale en octobre 2011, en publiant une annonce de recrutement qui révélait son adresse email personnelle rossulbricht@gmail.com. Cette erreur a permis aux autorités de l'identifier et contribué à son arrestation deux ans plus tard.}

\item \textbf{Question 2:} B
\\ \textit{Silk Road s'est démarqué des plateformes illégales existantes grâce à deux innovations techniques majeures : Tor, un réseau superposé confidentiel permettant de naviguer anonymement sur Internet, et Bitcoin, une cryptomonnaie décentralisée qui ne peut pas être arrêtée par les autorités. Les autres options sont incorrectes : PayPal et Western Union étaient utilisés par les anciennes plateformes comme Adamflowers mais n'offraient pas l'anonymat nécessaire ; Pecunix était initialement envisagé par Ross mais abandonné au profit de Bitcoin ; eBay et Amazon sont des plateformes légales qui ont inspiré le modèle de place de marché mais ne constituent pas des innovations techniques spécifiques à Silk Road.}

\item \textbf{Question 3:} D
\\ \textit{Ross Ulbricht avait créé un blog Wordpress à l'adresse silkroad420.wordpress.com pour fournir des instructions détaillées permettant d'accéder à la plateforme Silk Road via le réseau Tor. Ce blog a été rapidement suspendu par Wordpress, obligeant Ross à mettre en place un nouveau site web intermédiaire à l'adresse Silkroadmarket.org. Le chiffre 420 dans l'URL fait référence au code culturel associé au cannabis, ce qui était approprié pour une plateforme de vente de drogues.}

\item \textbf{Question 4:} C
\\ \textit{Ross Ulbricht gérait Good Wagon Books en 2010, une entreprise de revente de livres d'occasion qu'un ami lui avait cédée à Austin. Cet emploi alimentaire ne l'intéressait pas particulièrement, car il était bien plus obsédé par son projet personnel de création de Silk Road. Les autres noms proposés sont fictifs et ne correspondent à aucune entreprise mentionnée dans l'histoire de Ross Ulbricht.}

\item \textbf{Question 5:} B
\\ \textit{Samuel Edward Konkin III est le philosophe qui a théorisé l'agorisme dans les années 1970. Cette philosophie politique prône l'utilisation de l'économie souterraine comme moyen pacifique de réduire l'influence de l'État. Pour Ross Ulbricht, la découverte de l'agorisme fut une révélation car elle lui offrait un moyen concret d'agir selon ses convictions libertariennes. Bien que Ludwig von Mises, Murray Rothbard et Joseph Salerno soient tous des économistes de l'école autrichienne que Ross a étudiés, c'est spécifiquement Konkin III qui lui a fourni le cadre théorique pour créer Silk Road.}

\item \textbf{Question 6:} B
\\ \textit{Tor (acronyme de The Onion Router) était le réseau superposé confidentiel utilisé par Silk Road. Ce darknet basé sur le routage en oignon permettait de naviguer sur Internet et de communiquer de façon anonyme et relativement intraçable. Développé dans les années 2000, Tor était devenu assez populaire vers 2010 et se prêtait parfaitement au développement d'un marché noir en ligne. Les autres options sont incorrectes : I2P et Freenet sont d'autres réseaux anonymes mais n'étaient pas utilisés par Silk Road, tandis que VPN désigne un type de réseau privé virtuel différent.}

\item \textbf{Question 7:} C
\\ \textit{Silk Road était surnommée "l'Amazon de la drogue" en raison de son fonctionnement similaire à une place de marché classique comme Amazon, mais spécialisée dans la vente de drogues illicites. Ce surnom reflétait à la fois l'efficacité de la plateforme et sa spécialisation dans les substances interdites. Les autres options, bien que plausibles pour décrire une plateforme illégale, ne correspondent pas au surnom historiquement utilisé pour désigner Silk Road dans les médias et par les utilisateurs de l'époque.}

\item \textbf{Question 8:} C
\\ \textit{Ross Ulbricht avait initialement envisagé d'utiliser Pecunix, un système de devise en or numérique similaire à e-gold, pour son projet Silk Road en 2009. Cependant, il a découvert Bitcoin au cours de l'année 2010 et a réalisé que cette cryptomonnaie était mieux adaptée à son cas d'usage. Liberty Reserve était un autre système centralisé qui pouvait être arrêté par les autorités, e-gold était un prédécesseur de Pecunix, et PayPal était utilisé par d'autres plateformes rudimentaires mais ne convenait pas aux objectifs d'anonymat de Ross.}

\item \textbf{Question 9:} A
\\ \textit{Silk Road utilisait le service postal traditionnel pour expédier les produits, car la fouille systématique des colis postaux est interdite dans la plupart des pays. Les produits étaient camouflés astucieusement et les substances odorantes étaient placées dans des sachets sous vide pour éviter la détection par les chiens renifleurs. Cette méthode exploitait les protections légales existantes du courrier postal plutôt que de créer des systèmes de livraison alternatifs.}

\item \textbf{Question 10:} D
\\ \textit{Ross Ulbricht était profondément influencé par la philosophie agoriste de Samuel Edward Konkin III, qui prône l'utilisation de l'économie souterraine comme moyen pacifique de réduire l'influence de l'État. Dans ses écrits, il explique clairement que sa démarche était avant tout révolutionnaire et non opportuniste. Il voyait dans Silk Road un moyen d'appliquer concrètement les principes libertariens et agoristes, en créant un marché libre de toute intervention étatique. Les autres motivations proposées sont incorrectes car elles ne reflètent pas ses convictions philosophiques profondes.}

\item \textbf{Question 11:} C
\\ \textit{Le 11 octobre 2011, Ross Ulbricht a commis une erreur majeure en publiant une annonce de recrutement sur le forum Bitcoin sous le pseudonyme Altoid, le même qu'il avait utilisé pour promouvoir Silk Road. Dans cette annonce, il cherchait un développeur pour une startup Bitcoin et a inclus son adresse email personnelle rossulbricht@gmail.com, révélant ainsi son nom civil. Cette connexion directe entre le pseudonyme Altoid et son identité réelle a fourni aux enquêteurs un élément crucial pour l'identifier et l'arrêter deux ans plus tard.}

\item \textbf{Question 12:} B
\\ \textit{Black Market Reloaded était effectivement le concurrent de Silk Road lancé à la fin juin 2011, comme mentionné dans le texte. Underground Brokers était le nom initialement envisagé par Ross Ulbricht pour sa propre plateforme avant de choisir Silk Road. Adamflowers était une plateforme existante depuis 2006 mais qui utilisait des moyens de paiement classiques comme PayPal. Dark Web Market est un nom générique inventé qui n'apparaît pas dans l'histoire réelle de cette période.}

\item \textbf{Question 13:} B
\\ \textit{Mark Karpelès était un jeune développeur français né en 1985 qui avait déménagé au Japon en 2009 par attrait pour la culture locale. Il avait fondé la société Tibanne Co. Ltd. et gérait notamment un service d'hébergement appelé KalyHost. Son intérêt pour Bitcoin venait de sa curiosité technique plutôt que d'une motivation idéologique ou spéculative, contrairement à d'autres figures de l'époque.}

\item \textbf{Question 14:} C
\\ \textit{Le 19 juin 2011, un pirate a accédé au compte de Jed McCaleb sur Mt. Gox et a vendu massivement des bitcoins pour faire chuter le prix de 17,5 $ à 0,01 $. Cette manipulation lui a permis de retirer 2 000 BTC (soit 35 000 $ au cours de l'époque) en respectant la limite journalière de 1 000 $. Les autres montants correspondent à d'autres incidents : 25 000 BTC pour le vol d'Allinvain, 80 000 BTC pour le vol de mars 2011, et 17 000 BTC pour la perte de Bitomat.}

\item \textbf{Question 15:} B
\\ \textit{Contrairement à d'autres figures importantes de Bitcoin comme Martti Malmi ou Roger Ver qui étaient motivés par l'idéologie libertarienne, ou aux spéculateurs attirés par le profit, Mark Karpelès était principalement intéressé par les défis techniques. Il déclarait en 2014 que ce qui l'intéressait dans Bitcoin, c'étaient "les aspects techniques", notamment le maintien sécurisé d'une base de données mondiale, les portefeuilles privés sécurisés, et le système entièrement décentralisé. Cette curiosité technique explique pourquoi il s'est rapidement impliqué dans le développement communautaire, créant Bitcoin.it et travaillant sur diverses implémentations logicielles.}

\item \textbf{Question 16:} A
\\ \textit{Bitcoin-Central était la première plateforme européenne à véritablement concurrencer Mt. Gox, lancée le 29 décembre 2010 par David François en France. Bien que Britcoin soit apparue plus tard en mars 2011, Bitcoin-Central avait déjà établi sa position sur le marché européen avec support de l'euro et des virements SEPA. Bitomat était focalisée sur le marché polonais, et VirWoX était principalement dédiée aux monnaies virtuelles de jeux vidéos avant d'ajouter Bitcoin.}

\item \textbf{Question 17:} A
\\ \textit{Le 1er mars 2011, Mt. Gox a subi un vol de 80 000 bitcoins suite à un accès malveillant au fichier wallet.dat du site. Cette perte s'est produite juste avant l'officialisation du transfert de propriété le 6 mars. Jed McCaleb a informé Mark Karpelès de cet incident le 3 mars en lui disant que 'quelque chose d'horrible est arrivé'. Les autres montants correspondent à différents vols : 25 000 bitcoins pour Allinvain en juin, 50 000 $ pour le retrait non autorisé via Liberty Reserve, et 17 000 bitcoins pour la perte de Bitomat en juillet.}

\item \textbf{Question 18:} A
\\ \textit{Mark Karpelès gérait KalyHost, un service d'hébergement web, par l'intermédiaire de sa société Tibanne Co. Ltd. avant de reprendre Mt. Gox. Bien qu'il ait créé le wiki Bitcoin.it et travaillé sur QBitcoin après sa découverte de Bitcoin, ces projets étaient des contributions à la communauté Bitcoin plutôt que des services commerciaux qu'il gérait. Tibanne Co. Ltd. était sa société de freelance, mais pas un service spécifique.}

\item \textbf{Question 19:} A
\\ \textit{Le prix du bitcoin a atteint brièvement les 32 dollars le 8 juin 2011, marquant le sommet de ce qui sera appelé plus tard « la grande bulle de 2011 » par Vitalik Buterin. Cette hausse spectaculaire avait commencé en avril après l'article de Forbes, s'était poursuivie en mai pour atteindre 10 dollars, puis l'article de Gawker avait déclenché l'épisode spéculatif final jusqu'à 32 dollars. Les autres montants proposés ne correspondent pas à ce pic historique de juin 2011.}

\item \textbf{Question 20:} A
\\ \textit{Mark Karpelès avait fondé la société Tibanne Co. Ltd. au Japon, qui portait le nom de son chat Tibane. C'est par l'intermédiaire de cette société qu'il travaillait en freelance et gérait notamment un service d'hébergement appelé KalyHost. Lorsqu'il a repris Mt. Gox en février 2011, c'est Tibanne Co. Ltd. qui a acquis la plateforme de change. Les autres options sont incorrectes : KalyHost était le nom du service d'hébergement mais pas de la société, Bitcoin Consultancy était l'organisation fondée par Amir Taaki au Royaume-Uni, et QBitcoin était le nom d'une implémentation logicielle que Mark avait commencé à développer sans l'achever.}

\item \textbf{Question 21:} B
\\ \textit{Britcoin était la plateforme de change lancée le 26 mars 2011 par Amir Taaki au Royaume-Uni, spécialement conçue pour le marché britannique et permettant l'échange entre la livre sterling et le bitcoin. Intersango était le nom du logiciel open source qui soutenait la plateforme, mais pas le nom initial de la plateforme elle-même. Bitcoin-Central était la plateforme française de David François, et TradeHill était la plateforme américaine lancée plus tard en juin 2011.}

\item \textbf{Question 22:} B
\\ \textit{Selon le texte, entre la reprise effective de Mark Karpelès en mars 2011 et la fin du mois de mai, le nombre de comptes d'utilisateurs de Mt. Gox est passé de 3 000 à 60 000. Cette croissance spectaculaire illustre l'explosion de l'intérêt pour Bitcoin pendant la première bulle spéculative de 2011. Cette augmentation massive du nombre d'utilisateurs s'est accompagnée d'une croissance similaire du volume d'échange journalier, qui est passé de quelques milliers de dollars à plusieurs centaines de milliers de dollars sur la même période.}

\item \textbf{Question 23:} D
\\ \textit{Mark Karpelès est effectivement apparu dans le documentaire français 'Suck my Geek' en 2007, qui était consacré aux passionnés d'imaginaire, de technologie et d'informatique. À cette époque, il était encore un jeune geek 'mal dans sa peau' selon le texte, bien avant de découvrir Bitcoin et de reprendre Mt. Gox. Cette apparition télévisuelle illustre son profil de passionné de technologie qui l'a naturellement mené vers Bitcoin quelques années plus tard. Les autres titres proposés sont fictifs et n'existent pas.}

\item \textbf{Question 24:} D
\\ \textit{ClearCoin était le service de dépôt fiduciaire géré par Gavin Andresen. Après le krach éclair du 19 juin 2011 sur Mt. Gox, Kevin Day (alias Toasty) avait acheté 259 685 bitcoins à 0,0101 $ chacun. Suspecté par Mark Karpelès d'agir en collusion avec le pirate, il a prouvé sa bonne foi en plaçant les bitcoins sur un compte séquestre auprès de ClearCoin. Les fonds ont été libérés un mois plus tard, démontrant l'utilité de ce service d'arbitrage dans l'écosystème Bitcoin naissant.}

\item \textbf{Question 25:} B
\\ \textit{Roger Ver a été surnommé « le Jésus de Bitcoin » (Bitcoin Jesus) en février 2012 par le tradeur canadien GoWest, en raison de son prosélytisme intense pour promouvoir Bitcoin. Ce surnom lui est resté bien qu'il ne l'ait jamais revendiqué. Ver était devenu l'un des promoteurs les plus passionnés de la cryptomonnaie, multipliant les interventions publiques, finançant des publicités radio et investissant massivement dans l'écosystème Bitcoin. Les autres surnoms proposés, bien que similaires dans l'esprit, ne correspondent pas à celui qui lui a été historiquement attribué.}

\item \textbf{Question 26:} B
\\ \textit{Le mouvement libertarien américain a particulièrement favorisé l'adoption de Bitcoin car il prônait à la fois la liberté économique et la liberté des mœurs, constituant une troisième voie politique. Les libertariens jugeaient que l'intervention de l'État n'était pas désirable, y compris dans la fonction monétaire, et voyaient dans Bitcoin une alternative à la monnaie étatique grâce à son caractère non censurable et non inflationniste. Contrairement aux autres mouvements politiques mentionnés, les libertariens étaient idéologiquement alignés avec les principes de décentralisation et de limitation du pouvoir gouvernemental que représentait Bitcoin.}

\item \textbf{Question 27:} C
\\ \textit{Free Talk Live est l'émission de radio libertarienne animée par Ian Freeman et Mark Edge qui a introduit Bitcoin à Roger Ver. Roger Ver était un auditeur fidèle de cette émission, et c'est en écoutant les mentions de Bitcoin qui y étaient faites, notamment après l'émission du 16 mars 2011 consacrée à Bitcoin et Silk Road, qu'il a découvert la cryptomonnaie. InfoWars est mentionnée comme utilisant le même réseau de diffusion mais n'a pas joué ce rôle spécifique. The Bitcoin Show était une autre émission mais animée par Bruce Wagner, et Liberty Radio Network n'est pas mentionné dans le contexte de la découverte de Bitcoin par Roger Ver.}

\item \textbf{Question 28:} A
\\ \textit{Le New Hampshire a été choisi en 2003 par le Free State Project comme destination pour regrouper les libertariens américains. Ce choix s'explique par la culture individualiste de cet État, symbolisée par sa devise officielle « Live Free or Die » (Vivre libre ou mourir). Le Colorado, le Texas et la Californie, bien qu'ayant des communautés libertariennes, n'ont pas été sélectionnés par ce mouvement de migration politique qui visait à concentrer 20 000 libertariens dans un petit État pour influencer le gouvernement local.}

\item \textbf{Question 29:} C
\\ \textit{L'école autrichienne est effectivement l'école de pensée économique privilégiée par les libertariens américains. Cette école se caractérise par sa méthode rationnelle fondée sur le comportement individuel plutôt que sociétal, et rejette largement l'intervention de l'État. Les libertariens adhèrent particulièrement à l'idée que la dégradation du système monétaire au XXe siècle, passant de l'étalon-or aux monnaies fiduciaires, a provoqué une explosion des dépenses publiques. L'école de Chicago, bien que libérale, accepte davantage l'intervention étatique. L'école keynésienne prône au contraire une intervention active de l'État dans l'économie. L'école monétariste, associée à Milton Friedman, se concentre sur le contrôle de la masse monétaire mais n'exclut pas le rôle de l'État.}

\item \textbf{Question 30:} A
\\ \textit{Le slogan « Nous sommes les 99 \% » était la revendication principale du mouvement Occupy Wall Street, en opposition à « la cupidité et la corruption des 1 \% ». Ce mouvement a adopté Bitcoin pour échapper à la censure financière après que des comptes PayPal aient été gelés. « Vivre libre ou mourir » est la devise du New Hampshire, « La séparation de la monnaie et de l'État » est une expression d'Erik Voorhees sur Bitcoin, et « Occupons la Réserve fédérale » n'est pas un slogan réel du mouvement.}

\item \textbf{Question 31:} D
\\ \textit{Les libertariens attachés à l'école autrichienne estimaient que Bitcoin violait le théorème de régression de Ludwig von Mises, car il n'avait pas de valeur intrinsèque ni d'adossement à quelque chose de concret comme l'or. Ils pensaient qu'une vraie monnaie devait avoir eu une utilité matérielle avant de devenir un moyen d'échange. Les autres options sont incorrectes car elles ne correspondent pas aux critiques théoriques formulées par les partisans de l'école autrichienne.}

\item \textbf{Question 32:} D
\\ \textit{Le mouvement Occupy Wall Street a été créé en réaction aux désastres de la récession économique provoquée par la crise des subprimes. Face à cette situation, la fondation Adbusters a appelé à occuper Wall Street pour protester contre l'influence des grandes entreprises et des puissances financières jugées responsables de la crise. Les autres réponses correspondent à des événements liés à Bitcoin mais qui ne sont pas à l'origine du mouvement Occupy Wall Street.}

\item \textbf{Question 33:} D
\\ \textit{FeedZeBirds était une plateforme publicitaire lancée le 23 novembre 2011 par Erik Voorhees et Ira Miller. Cette plateforme permettait de payer en Bitcoin pour mettre en avant un contenu sur Twitter en achetant des retweets. BitMunchies était un magasin en ligne de nourriture créé par Ira Miller en mai 2011, BTCinch était un processeur de paiement développé pendant l'été 2011, et Coinapult était un service d'envoi de fonds par e-mail développé plus tard par les deux associés.}

\item \textbf{Question 34:} B
\\ \textit{Ron Paul était surnommé « Doctor No » parce qu'il avait l'habitude de voter contre tous les projets de loi qui réduisaient la liberté des citoyens ou qui augmentaient les dépenses publiques. Ce surnom reflétait sa position systématiquement opposée aux mesures gouvernementales expansives, cohérente avec ses convictions libertariennes. Sa campagne présidentielle de 2012 a coïncidé avec la période d'essor de Bitcoin et a inspiré de nombreux libertariens qui ont ensuite adopté la cryptomonnaie.}

\item \textbf{Question 35:} D
\\ \textit{Le New Hampshire a été choisi par le Free State Project en 2003 en raison de sa culture individualiste, symbolisée par sa devise officielle "Live Free or Die" (Vivre libre ou mourir). Cette devise reflétait parfaitement les valeurs libertariennes du mouvement qui cherchait un État où regrouper les partisans de la liberté. Les autres devises proposées, bien qu'ayant des connotations libertariennes, ne correspondent pas à la devise officielle du New Hampshire et n'ont donc pas influencé ce choix stratégique.}

\item \textbf{Question 36:} A
\\ \textit{Les libertariens traditionnels, attachés à l'école autrichienne, rejetaient Bitcoin car il contredisait le théorème de régression de Ludwig von Mises. Selon eux, une monnaie doit avoir une valeur intrinsèque et être adossée à quelque chose de concret comme l'or ou l'argent. Bitcoin, étant purement numérique sans valeur matérielle antérieure, ne pouvait pas devenir une vraie monnaie selon cette théorie. Les autres options sont incorrectes car Bitcoin est décentralisé, non-inflationniste et indépendant de l'État.}

\item \textbf{Question 37:} D
\\ \textit{Le chapitre indique clairement qu'en juin 2011, le piratage de Mt. Gox a stoppé net l'engouement spéculatif, puis que le fiasco lié à la fermeture de MyBitcoin fin juillet a finalement déclenché le dégonflement de la bulle et conduit au premier réel marché baissier. Ces deux événements consécutifs ont causé la chute du prix de 32$ à 6$ puis jusqu'à 2$ en novembre. Les autres réponses mentionnent des événements qui se sont produits après le début du marché baissier ou qui ne correspondent pas aux faits historiques présentés.}

\item \textbf{Question 38:} B
\\ \textit{D'après le texte, le Bitcoin a atteint un sommet de 32 dollars le 8 juin 2011, puis a chuté drastiquement suite au piratage de Mt. Gox et à la fermeture de MyBitcoin. Le prix est d'abord descendu à 6 dollars début août, puis a continué à baisser pour finalement atteindre environ 2 dollars en novembre 2011. Cette chute représentait effectivement une baisse de plus de 90\% par rapport au pic de juin, illustrant l'ampleur du premier véritable marché baissier de Bitcoin.}

\item \textbf{Question 39:} B
\\ \textit{BitInstant était le service de change instantané développé par le Gallois Gareth Nelson et l'Américain Charlie Shrem en 2011. Ce service résolvait le problème des virements internationaux lents vers Mt. Gox en servant d'intermédiaire qui virait les fonds sur les plateformes de change une fois reçus dans le pays d'origine. BitInstant a ouvert ses portes au public le 23 août 2011 et est rapidement devenu populaire, transférant 1,2 million de dollars dès janvier 2012.}

\item \textbf{Question 40:} A
\\ \textit{Zhou Tong, de son vrai nom Ryan Tong Zhou, était un surdoué chinois de 17 ans qui étudiait à Singapour. Il a lancé Bitcoinica le 8 septembre 2011, une plateforme révolutionnaire qui offrait pour la première fois le trading sur marge avec effet de levier et la vente à découvert dans l'écosystème Bitcoin. James McCarthy (Nefario) était le créateur de GLBSE, Mircea Popescu a fondé MPEx, et Charlie Shrem était le cofondateur de BitInstant.}

\item \textbf{Question 41:} D
\\ \textit{GLBSE, acronyme de GLobal Bitcoin Stock Exchange, était la plateforme créée par James McCarthy (pseudonyme Nefario) en avril 2011. Cette plateforme permettait d'émettre des actions pour lever des capitaux et de verser des dividendes aux actionnaires. MPEx était une plateforme similaire mais créée par Mircea Popescu. BitInstant était un service de change instantané créé par Gareth Nelson et Charlie Shrem. Bitcoinica était une plateforme de trading sur marge créée par Zhou Tong.}

\item \textbf{Question 42:} C
\\ \textit{Le Bitcoin a atteint son sommet à 32 dollars le 8 juin 2011 avant que le piratage de Mt. Gox et la fermeture de MyBitcoin ne déclenchent l'éclatement de la bulle. Ce pic de 32 dollars était remarquable car le prix a ensuite chuté drastiquement jusqu'à 6 dollars début août, puis jusqu'à 2 dollars en novembre, représentant une baisse de plus de 90\% par rapport à ce sommet. Les autres prix proposés (16, 48 et 64 dollars) ne correspondent pas au sommet historique de cette première grande bulle spéculative de Bitcoin.}

\item \textbf{Question 43:} D
\\ \textit{BitInstant a été créé spécifiquement pour résoudre le problème des délais de transfert d'argent vers Mt. Gox. Comme cette plateforme était installée au Japon, les clients occidentaux devaient effectuer des virements internationaux qui demandaient plusieurs jours d'attente, ce qui était particulièrement frustrant dans un marché aussi volatil que celui du Bitcoin. BitInstant servait d'intermédiaire en recevant les fonds localement puis en les virant rapidement sur Mt. Gox, permettant ainsi un accès quasi-instantané aux dollars sur la plateforme de change japonaise.}

\item \textbf{Question 44:} B
\\ \textit{L'expression 'get Zhou Tonged' (se faire zhou-tonguer) est devenue populaire pour décrire les liquidations sur Bitcoinica, en référence au nom de son créateur Zhou Tong. Cette expression était si connue qu'elle a même donné son nom à une chaîne YouTube lancée en mars 2012. Les autres expressions proposées n'existaient pas à cette époque, même si 'se faire mt-goxer' deviendra plus tard une référence aux problèmes de Mt. Gox.}

\item \textbf{Question 45:} A
\\ \textit{GLBSE et MPEx ont été créées spécifiquement pour libéraliser l'investissement en permettant aux entreprises de l'écosystème Bitcoin d'émettre des actions et des obligations sans passer par les autorités de régulation traditionnelles comme la SEC. Ces plateformes exploitaient le flou juridique autour de Bitcoin pour créer un modèle financier alternatif selon l'approche agoriste. Les autres options décrivent des services différents : les échanges de devises étaient assurés par des plateformes comme Mt. Gox, le trading sur marge par Bitcoinica, et les portefeuilles par d'autres développements techniques.}

\item \textbf{Question 46:} C
\\ \textit{Bitcoinica était révolutionnaire car elle était la première plateforme à offrir du trading sur marge dans l'écosystème Bitcoin, permettant aux utilisateurs d'utiliser un effet de levier pour amplifier leurs gains (et leurs pertes) et d'effectuer de la vente à découvert. Cette innovation a rendu possible la spéculation avancée et la couverture de risque, contrairement aux plateformes existantes qui ne proposaient que de simples échanges. Les autres options correspondent respectivement à BitInstant, GLBSE/MPEx, et encore BitInstant.}

\item \textbf{Question 47:} B
\\ \textit{Paul Krugman, économiste néokeynésien et prix Nobel, a critiqué Bitcoin en établissant un parallèle avec l'étalon-or. Il arguait que Bitcoin avait 'créé son propre modèle privé d'étalon-or, dans lequel la masse monétaire est fixe' et que cette limitation encourageait la thésaurisation plutôt que l'utilisation comme moyen d'échange. Il faisait référence à l'ouvrage 'Golden Fetters' de Barry Eichengreen qui soutenait que l'étalon-or avait contribué à la Grande Dépression des années 1930. Les autres réponses, bien que représentant des critiques communes de Bitcoin, ne correspondent pas à l'argumentation spécifique de Krugman dans cet article.}

\item \textbf{Question 48:} A
\\ \textit{Mircea Popescu a financé le développement de MPEx par une méthode innovante pour l'époque : une levée de fonds en émettant des actions et des obligations libellées en bitcoins. Cette opération lui a permis de recueillir plus de 3 200 bitcoins (environ 16 000 $ à l'époque) en février 2012. Cette approche illustrait parfaitement la philosophie de contournement des régulations financières traditionnelles prônée par les entrepreneurs Bitcoin de l'époque, utilisant la cryptomonnaie pour créer un système financier alternatif.}

\item \textbf{Question 49:} D
\\ \textit{La vérification de paiement simplifiée (SPV) était une innovation cruciale décrite dans le livre blanc de Bitcoin qui permettait aux clients légers de vérifier que leurs transactions appartenaient bien à la chaîne de blocs sans avoir besoin de télécharger et stocker l'intégralité des données. Cette technologie, implémentée dans BitCoinJ par Mike Hearn, était essentielle pour les appareils mobiles avec des ressources limitées. Les autres options décrivent des fonctionnalités différentes : le chiffrement AES concernait les portefeuilles web, les adresses personnalisées étaient générées par VanityGen, et le mélange était fourni par des services spécialisés comme BitLaundry.}

\item \textbf{Question 50:} D
\\ \textit{Les portefeuilles déterministes ont été développés pour résoudre le problème des sauvegardes incomplètes. Avec le système traditionnel, chaque nouvelle clé était générée aléatoirement, obligeant l'utilisateur à sauvegarder régulièrement son fichier wallet.dat. Les portefeuilles déterministes permettent de dériver toutes les clés à partir d'une seule graine, nécessitant ainsi une seule sauvegarde initiale pour récupérer l'intégralité des fonds.}

\item \textbf{Question 51:} B
\\ \textit{BCCAPI souffrait de deux défauts majeurs selon le texte : premièrement, le code de l'infrastructure logicielle du serveur n'était pas public, ce qui empêchait le déploiement d'autres serveurs par la communauté. Deuxièmement, le système conservait en mémoire les transactions liées aux adresses des portefeuilles pour fournir le solde total, ce qui posait un problème de confidentialité. Ces limitations ont poussé Thomas Voegtlin à créer Electrum avec un réseau de serveurs plus ouvert et respectueux de la vie privée.}

\item \textbf{Question 52:} C
\\ \textit{Les portefeuilles déterministes révolutionnent la gestion des clés en dérivant toutes les clés privées de manière déterministe à partir d'une information unique appelée graine (seed), contrairement aux portefeuilles traditionnels qui génèrent chaque clé de façon aléatoire et indépendante. Cette approche simplifie drastiquement la sauvegarde puisqu'il suffit de conserver la graine pour récupérer l'intégralité des fonds, éliminant le besoin de sauvegarder constamment le fichier wallet.dat. Cette innovation répond au dilemme sécurité-accessibilité en réduisant les risques de perte de fonds dus à des sauvegardes insuffisantes.}

\item \textbf{Question 53:} B
\\ \textit{Instawallet se distinguait par son système d'accès simplifié : contrairement aux autres services qui demandaient une inscription avec identifiant et mot de passe, Instawallet générait automatiquement une URL unique lors de la première visite, comme 'https://instawallet.org/w/rq2SB02ai6BnWaEBywAlP52cw7qwUAA', qui était sauvegardée dans le navigateur. Cette approche sans inscription était révolutionnaire pour l'époque et rendait le service extrêmement facile d'utilisation.}

\item \textbf{Question 54:} B
\\ \textit{Les bitcoins physiques de Casascius utilisaient un système d'hologramme personnalisé qui recouvrait la clé privée. Cet hologramme était conçu pour s'endommager de manière visible dès qu'on tentait de l'enlever ou de le manipuler pour accéder à la clé privée en dessous. Cette méthode permettait de détecter immédiatement si quelqu'un avait tenté d'accéder aux fonds stockés sur la pièce, garantissant ainsi l'intégrité du bitcoin physique. Les autres options mentionnent des méthodes qui n'étaient pas utilisées par Casascius : pas de chiffrement par mot de passe, pas de serveur central requis, et pas de technologie d'encre invisible.}

\item \textbf{Question 55:} A
\\ \textit{Gregory Maxwell a conçu deux types de dérivation pour les portefeuilles déterministes. Le type 1 générait directement les clés privées à partir de la graine, tandis que le type 2 était plus sophistiqué car il permettait de dériver les adresses publiques à partir d'une clé publique maîtresse obtenue depuis la graine, sans jamais exposer la graine elle-même. Cette innovation était particulièrement utile pour les logiciels de traitement de paiements qui pouvaient générer une nouvelle adresse à chaque transaction sans que les clés privées ne soient présentes sur le système, améliorant considérablement la sécurité.}

\item \textbf{Question 56:} D
\\ \textit{Bitcoin Fog se distinguait de BitLaundry par son approche du mélange de bitcoins. Contrairement à BitLaundry qui utilisait des adresses à usage unique pour un mélange immédiat, Bitcoin Fog imposait aux utilisateurs de maintenir un solde sur leur compte. Cette conception permettait d'effectuer des retraits différés et échelonnés sur une période de 6 à 96 heures, vers plusieurs adresses différentes. Cette méthode rendait le traçage des fonds beaucoup plus difficile car elle introduisait une dimension temporelle et multipliait les points de sortie, améliorant ainsi l'efficacité du mélange par rapport aux solutions existantes.}

\item \textbf{Question 57:} C
\\ \textit{StrongCoin était le premier portefeuille web à exploiter la bibliothèque BitcoinJS, mais son principal défaut était sa structure tarifaire excessive. Le service prélevait des frais de 1\% sur chaque montant envoyé, avec un plafond pouvant atteindre 1 BTC par transaction. Cette commission élevée constituait un obstacle majeur à l'adoption, contrairement à My Wallet de Blockchain.info qui ne prélevait aucune commission. Bien que les deux services offraient des fonctionnalités similaires de chiffrement des clés privées côté client, la gratuité de My Wallet combinée à la clarté de l'explorateur de blocs intégré explique son succès rapide avec 5000 comptes dès mars 2012.}

\item \textbf{Question 58:} B
\\ \textit{Ben Reeves a développé Blockchain.info en août 2011 avec des fonctionnalités spécifiques que les autres explorateurs n'offraient pas : l'inclusion des blocs orphelins permettant de détecter les tentatives de double dépense, et l'estimation du volume réel de bitcoins transactés plutôt que simplement le montant envoyé. Ces innovations techniques lui ont permis de concurrencer efficacement le Bitcoin Block Explorer (BBE) de Theymos et ABE, devenant rapidement l'explorateur préféré de la communauté avant d'intégrer la fonctionnalité de portefeuille My Wallet en décembre.}

\item \textbf{Question 59:} B
\\ \textit{BCCAPI était révolutionnaire car elle intégrait la génération déterministe des clés privées, permettant de reconstituer toutes les clés à partir d'une seule information initiale. Cette approche éliminait le besoin de faire des sauvegardes répétées du fichier wallet.dat, contrairement aux portefeuilles traditionnels qui généraient les clés aléatoirement. Les autres options décrivent des fonctionnalités qui existaient dans d'autres solutions mais n'étaient pas la spécificité principale de BCCAPI.}

\item \textbf{Question 60:} C
\\ \textit{Mike Caldwell a explicitement déclaré en 2013 que son objectif était d'utiliser les Casascius Coins comme outil pédagogique pour aider le monde à visualiser une pièce de monnaie virtuelle avec un objet tangible fonctionnel. Bien que les bitcoins physiques offraient certains avantages de sécurité et de praticité, leur fonction première était éducative et démonstrative, permettant d'attirer un public différent comme les amateurs de numismatique et de rendre Bitcoin plus compréhensible au grand public.}

\item \textbf{Question 61:} A
\\ \textit{Le modèle proportionnel simple de Slush récompensait les utilisateurs au prorata des parts trouvées depuis le dernier bloc, mais cette méthode permettait aux mineurs de "tricher" par pool hopping. Ils pouvaient se déconnecter après avoir soumis un certain nombre de parts et rediriger leur puissance vers leur propre compte ou une autre coopérative, obtenant ainsi environ 28\% de bitcoins supplémentaires. C'est pourquoi Slush a dû implémenter un système de score "anti-tricherie" le 4 février 2011, donnant plus de poids aux parts récentes.}

\item \textbf{Question 62:} C
\\ \textit{Le protocole Stratum a été proposé par Slush le 11 septembre 2012 comme une alternative plus légère et efficace aux solutions existantes. Contrairement à getblocktemplate (proposé par Luke-Jr) ou getmemorypool, Stratum permettait une communication plus efficace en termes de bande passante entre les clients légers des participants et les nœuds des coopératives. Le long polling était une technique antérieure utilisée avec getwork, mais n'était pas un protocole complet. Stratum s'est rapidement imposé, représentant 74\% du taux de hachage en décembre 2012.}

\item \textbf{Question 63:} D
\\ \textit{Luke-Jr a développé le modèle SMPPS (Shared Maximum Pay Per Share) pour Eligius en juin 2011, après avoir expérimenté avec le MPPS (Maximum Pay Per Share) qui s'était révélé défectueux. Le SMPPS utilise un tampon collectif pour ne jamais distribuer plus que ce que la coopérative possède réellement, évitant ainsi les risques financiers du modèle PPS classique tout en offrant une rémunération plus stable que le modèle proportionnel. Les autres options mentionnées sont des méthodes développées par d'autres coopératives : le PPLNS par Mineco.in, le CPPSRB par BitPenny, et la méthode géométrique double par EclipseMC.}

\item \textbf{Question 64:} A
\\ \textit{Eligius a innové en inscrivant son nom dans le champ d'entrée de la transaction de récompense du bloc 130 635, pratique appelée 'signature de la coinbase'. Cette méthode permettait de rendre le fonctionnement des coopératives plus transparent en identifiant clairement quel groupe avait miné chaque bloc. Les autres options correspondent à des innovations différentes : Stratum fut développé par Slush en 2012, le PPLNS par Wuked de Mineco.in, et la sharechain par P2Pool de Forrest Voight.}

\item \textbf{Question 65:} D
\\ \textit{La coopérative Deepbit de Tycho a atteint son pic de domination le 5 août 2011 avec 52,86\% de la puissance de calcul mesurée du réseau Bitcoin. Cette concentration exceptionnelle a provoqué des inquiétudes dans la communauté concernant la centralisation du minage et la sécurité du réseau. Bien que certains aient craint une compromission de la sécurité de Bitcoin, les utilisateurs plus expérimentés ont rappelé que l'opérateur n'avait aucun intérêt économique à attaquer le réseau et que les participants restaient libres de changer de coopérative à tout moment.}

\item \textbf{Question 66:} B
\\ \textit{P2Pool révolutionnait le concept de coopérative de minage en éliminant le point de défaillance unique que représentait l'opérateur central. Contrairement aux coopératives traditionnelles où un seul opérateur choisit les transactions et construit les blocs, P2Pool utilisait une "sharechain" distribuée maintenue par tous les participants connectés en pair à pair. Ce système décentralisé empêchait qu'une seule entité contrôle la sélection des transactions et la construction des blocs, répondant ainsi aux préoccupations de centralisation qui inquiétaient la communauté Bitcoin, notamment après la domination de Deepbit qui avait atteint plus de 50\% du taux de hachage.}

\item \textbf{Question 67:} A
\\ \textit{La méthode PPLNS était innovante car elle ne se basait plus sur le cycle de production des blocs pour calculer les récompenses, mais sur une fenêtre fixe des N dernières parts soumises à la coopérative. Cette approche permettait de maintenir les avantages du modèle proportionnel tout en réduisant la vulnérabilité au pool hopping, contrairement au Pay Per Share qui payait immédiatement, à la méthode géométrique qui utilisait un score décroissant, ou au SMPPS qui employait un tampon collectif.}

\item \textbf{Question 68:} B
\\ \textit{P2Pool révolutionnait l'architecture des coopératives en remplaçant le modèle client-serveur traditionnel par un système décentralisé. Sa sharechain, similaire à la blockchain Bitcoin mais avec des blocs toutes les 5 secondes, était maintenue collectivement par tous les participants connectés en réseau pair à pair. Cette architecture éliminait le point de défaillance unique des coopératives centralisées et permettait aux mineurs de garder le contrôle total sur leur activité, contrairement aux autres systèmes où un opérateur central gérait la construction des blocs.}

\item \textbf{Question 69:} B
\\ \textit{Avec la généralisation du minage par GPU et l'augmentation exponentielle de la difficulté du réseau, les petits mineurs individuels faisaient face à une variance trop importante : ils pouvaient attendre des mois sans jamais trouver de bloc et donc sans recevoir de récompense. Les coopératives permettaient de lisser ces revenus en mutualisant la puissance de calcul et en redistribuant les gains proportionnellement aux contributions, garantissant ainsi des paiements plus réguliers et prévisibles.}

\item \textbf{Question 70:} D
\\ \textit{La sharechain de P2Pool était configurée avec une difficulté initialement 600 fois inférieure à celle du réseau Bitcoin principal. Cette différence permettait aux mineurs de soumettre des preuves de travail plus fréquemment (toutes les 5 secondes en moyenne) tout en maintenant la cohérence du système décentralisé. Cette configuration était essentielle pour assurer le bon fonctionnement du protocole pair-à-pair de P2Pool, où chaque part devait contenir l'empreinte de la part précédente pour maintenir l'intégrité de la chaîne.}

\item \textbf{Question 71:} C
\\ \textit{Luke-Jr a cessé d'inclure des prières dans les blocs d'Eligius suite aux critiques de membres de la communauté, notamment Graeme Tee d'Ozcoin, qui estimaient que cette pratique pouvait rebuter certaines personnes et que les mineurs participants n'étaient pas informés de ces ajouts. Bien que Luke-Jr ait d'abord défendu sa position en affirmant que les catholiques ne croient pas en la liberté de religion, il a finalement tenu compte des remarques et a donné aux mineurs la possibilité de demander l'inclusion de leurs propres messages.}

\item \textbf{Question 72:} B
\\ \textit{Le système P2Pool utilisait une sharechain avec un intervalle moyen de 5 secondes entre les parts, soit 120 fois plus rapide que l'intervalle de 10 minutes des blocs Bitcoin. Cette fréquence élevée permettait une meilleure répartition des récompenses et réduisait la variance pour les mineurs participants. La difficulté de la sharechain était ajustée régulièrement pour maintenir cet intervalle constant, contrairement aux autres propositions qui utilisaient des intervalles plus longs et moins optimisés pour le minage décentralisé.}

\item \textbf{Question 73:} A
\\ \textit{Russell O'Connor a identifié deux défauts majeurs dans OP_EVAL : premièrement, le risque de récursion qui pourrait faire tourner le programme indéfiniment (il a d'ailleurs trouvé une faille de récursion en seulement 70 minutes d'examen), et deuxièmement, l'impossibilité d'analyser statiquement le script pour compter les opérateurs de signature avant son exécution. Ces limitations rendaient OP_EVAL dangereux et imprévisible, contrairement aux autres options proposées qui étaient rétrocompatibles et n'augmentaient pas significativement la taille des transactions.}

\item \textbf{Question 74:} A
\\ \textit{La principale critique de Russell O'Connor concernant OP_EVAL était qu'il rendait impossible l'analyse statique des scripts. Avec OP_EVAL, il n'y avait aucun moyen d'analyser le script (par exemple pour compter le nombre d'utilisations d'OP_CHECKSIG) sans l'exécuter réellement. P2SH résolvait ce problème en permettant une analyse préalable du script. Les autres réponses sont incorrectes : P2SH n'était pas plus simple à implémenter, n'avait pas le soutien unanime (Luke-Jr s'y opposait fortement), et les deux approches étaient des soft forks évitant la scission de chaîne.}

\item \textbf{Question 75:} D
\\ \textit{Russell O'Connor a identifié que l'opérateur OP_EVAL empêchait l'analyse statique des scripts, car il fallait exécuter le script pour connaître son contenu réel. Cela posait un problème majeur pour compter les opérateurs de signature coûteux en ressources avant l'exécution. En revanche, Pay to Script Hash permettait cette analyse préalable en reconnaissant une forme spécifique de script. Les autres options sont incorrectes : OP_EVAL était rétrocompatible via les opcodes OP_NOP, il n'y avait pas de limitation stricte à 3 participants, et la taille des adresses n'était pas le problème principal soulevé.}

\item \textbf{Question 76:} C
\\ \textit{L'activation de P2SH nécessitait un signalement majoritaire des mineurs pour éviter une scission de chaîne. Tycho, opérateur de DeepBit (la plus grande coopérative avec ~30\% de la puissance), refusait initialement de déployer P2SH, préférant attendre un consensus plus large. Sans son soutien, le seuil de 55\% requis ne pouvait être atteint. Ce n'est qu'en mars 2012, après le déploiement du BIP 30, que Tycho a finalement accepté d'implémenter P2SH, permettant l'activation le 1er avril 2012.}

\item \textbf{Question 77:} A
\\ \textit{P2SH introduisait un concept révolutionnaire : au lieu d'exposer la complexité du script dans l'adresse de réception, seule l'empreinte (hash) du script était utilisée pour créer une adresse courte. Le script complet n'était révélé qu'au moment de la dépense, permettant des transactions complexes avec des adresses simples et courtes, similaires aux adresses Bitcoin classiques. Cette approche contrastait avec les multisignatures brutes où toute la complexité était visible dès la création de l'adresse.}

\item \textbf{Question 78:} B
\\ \textit{L'activation de P2SH utilisait un mécanisme de signalement par les mineurs, où ils devaient inclure la chaîne '/P2SH/' dans les blocs qu'ils produisaient. Gavin Andresen avait fixé un seuil de 55\% de signalement nécessaire pour déclencher l'activation. Ce mécanisme permettait d'assurer qu'une majorité de la puissance de calcul supportait le changement avant son activation, évitant ainsi une scission de chaîne. Les autres options sont incorrectes car le vote communautaire n'était que consultatif, il n'y avait pas d'activation automatique sans signalement, et l'unanimité des développeurs n'était pas requise ni obtenue.}

\item \textbf{Question 79:} D
\\ \textit{Le BIP 30 était une correction d'urgence pour résoudre le problème des transactions en double qui existaient sur la blockchain Bitcoin. Russell O'Connor avait découvert une vulnérabilité permettant de supprimer des transactions en exploitant ces doublons d'identifiants. Le BIP 30 interdisait donc la création future de transactions ayant le même identifiant qu'une transaction précédente non entièrement dépensée, empêchant ainsi cette attaque. Ce correctif était prioritaire et a temporairement interrompu le déploiement de Pay to Script Hash.}

\item \textbf{Question 80:} A
\\ \textit{Les opérateurs OP_NOP étaient des instructions muettes ajoutées par Satoshi qui n'invalidaient pas les scripts. Gavin a découvert qu'en transformant OP_NOP1 en OP_EVAL, puis plus tard en utilisant ce principe pour P2SH, il était possible de faire des modifications 'rétrocompatibles' (soft forks) où seule une chaîne survivrait si la puissance de calcul majoritaire appliquait le changement. Cela évitait les scissions de chaîne dangereuses qui auraient pu diviser le réseau Bitcoin.}

\item \textbf{Question 81:} C
\\ \textit{Après le départ de Satoshi, le développement de Bitcoin est devenu un effort collaboratif. Alors que Satoshi était le seul développeur principal, Gavin Andresen a pris la relève en tant que mainteneur principal mais s'est entouré d'une équipe de développeurs 'attitrés' incluant Pieter Wuille, Nils Schneider, Jeff Garzik et Wladimir van der Laan. Ces programmeurs bénéficiaient d'un accès en écriture au dépôt GitHub et représentaient Bitcoin au niveau technique. Cette transition vers un modèle collaboratif a d'ailleurs contribué aux tensions lors de la bataille pour P2SH, car il y avait désormais assez de développeurs pour avoir des opinions divergentes sur l'évolution du protocole.}

\item \textbf{Question 82:} A
\\ \textit{Tycho justifiait son opposition à P2SH en donnant trois raisons principales : l'absence de consensus entre les développeurs ou entre les mineurs sur la manière d'implémenter le paiement aux scripts, le caractère quelque peu 'bricolé' de la solution proposée, et son refus de devenir la seule entité à décider de cette question. Il ne s'agissait pas d'une volonté de domination du réseau ou de considérations économiques, mais bien d'une position technique et de gouvernance. Tycho avait d'ailleurs soutenu la proposition originale d'OP_EVAL mais n'approuvait ni le BIP 16 ni le BIP 17.}

\item \textbf{Question 83:} D
\\ \textit{Bitcoin-Qt, développé par Wladimir van der Laan à partir de juin 2011, constituait une refonte complète de l'interface graphique de Bitcoin. La principale innovation était le remplacement de la bibliothèque wxWindows, utilisée par Satoshi dans les versions originales, par le framework graphique Qt. Cette transition technique a permis d'améliorer significativement l'expérience utilisateur et la portabilité du logiciel. L'interface Qt est devenue officielle dans la version 0.5 sortie le 21 novembre 2011, marquant une étape importante dans l'évolution du logiciel Bitcoin vers une interface plus moderne et professionnelle.}

\item \textbf{Question 84:} B
\\ \textit{Russell O'Connor a découvert qu'il était possible d'exploiter les transactions de récompense en double présentes sur la blockchain pour supprimer complètement toutes les transactions qui en découlaient, en utilisant une réorganisation de chaîne. Cette vulnérabilité critique a nécessité la création urgente du BIP 30 pour interdire aux blocs futurs de contenir des transactions avec des identifiants correspondant à des transactions antérieures non entièrement dépensées. Les autres réponses concernent des problèmes différents : le bug d'OP_CHECKMULTISIG était connu mais non critique, Bitcoin-Qt était une interface graphique sans faille de sécurité majeure, et le signalement des mineurs ne présentait pas de vulnérabilité de falsification.}

\item \textbf{Question 85:} D
\\ \textit{Namecoin a introduit le minage combiné (merged mining), un procédé proposé par Satoshi Nakamoto lui-même qui permet de réutiliser le travail de minage effectué sur Bitcoin pour valider simultanément les blocs de chaînes auxiliaires plus petites. Ce système a été activé sur Namecoin le 8 octobre 2011. La preuve d'enjeu hybride était utilisée par PPCoin, l'algorithme Scrypt par Tenebrix et Litecoin, et les nœuds de confiance par SolidCoin 2.0.}

\item \textbf{Question 86:} D
\\ \textit{Tenebrix a été la première cryptomonnaie à implémenter la fonction de hachage Scrypt, développée par Colin Percival pour Tarsnap. Cette fonction, grâce à sa forte demande de mémoire, réduisait drastiquement l'efficacité des processeurs graphiques (GPU) dans le minage, permettant ainsi aux processeurs centraux (CPU) d'être utilisés de manière profitable. Les autres options correspondent à des innovations d'autres cryptomonnaies : la preuve d'enjeu hybride était l'innovation de PPCoin, le temps de bloc de 15 secondes était caractéristique de GeistGeld, et les nœuds de confiance étaient une particularité de SolidCoin 2.0.}

\item \textbf{Question 87:} C
\\ \textit{La différence fondamentale entre ces deux cryptomonnaies était le préminage. Thomas Nasakioto avait secrètement miné les 6 055 premiers blocs d'Ixcoin pendant trois mois avant son lancement public, générant plus de 580 000 ixcoins. Cette pratique controversée a motivé le développement d'I0coin comme une alternative équitable, explicitement présentée comme un fork d'Ixcoin 'avec 0 bloc prégénéré'. Les autres options sont incorrectes : les deux utilisaient SHA-256, I0coin avait un temps de bloc de 5 minutes (pas 2,5), et aucune n'était compatible avec le minage combiné à l'époque.}

\item \textbf{Question 88:} A
\\ \textit{Le lancement de Litecoin a été marqué par un minage accéléré problématique : plus de 10 000 blocs ont été extraits en moins de 24 heures au lieu des 12,5 jours prévus, générant 500 000 litecoins soit 0,6\% de la quantité totale planifiée. Ce problème était dû à une difficulté de minage initiale trop basse combinée à un algorithme d'ajustement limité qui ne pouvait multiplier la difficulté que par 4 maximum. De plus, le bloc de genèse était horodaté au 9 octobre, causant un premier ajustement à la baisse plutôt qu'à la hausse.}

\item \textbf{Question 89:} C
\\ \textit{Dans sa citation de mars 2013, Sunny King identifie spécifiquement "l'introduction d'un registre public décentralisé" comme la principale avancée de Bitcoin. Bien que la preuve de travail, SHA-256 et le réseau pair-à-pair soient des composants techniques importants de Bitcoin, c'est le concept de registre décentralisé qui constitue l'innovation fondamentale permettant de créer des monnaies virtuelles résistantes aux tentatives de fermeture, selon l'analyse de Sunny King.}

\item \textbf{Question 90:} A
\\ \textit{L'expression "yet another cryptocurrency" (une cryptomonnaie de plus) est apparue dans la communauté pour exprimer l'exaspération face à la multiplication des cryptomonnaies alternatives. Cette expression reprend la formule "yet another" déjà répandue dans le monde du logiciel libre pour décrire "quelque chose qui existe déjà en trop grand nombre". Elle a été utilisée de manière péjorative contre diverses cryptomonnaies comme Namecoin, Tenebrix, ou Liquidcoin, et a même inspiré le nom d'une cryptomonnaie appelée YACoin ("yet another coin") lancée en 2013.}

\item \textbf{Question 91:} D
\\ \textit{Le triangle de Zooko énonce qu'un système d'identifiants ne peut simultanément être sécurisé (un nom correspond à un seul élément), décentralisé (pas d'autorité centrale nécessaire) et intelligible pour l'homme (noms mémorisables). Bitcoin a permis de résoudre ce trilemme, rendant possible la "quadrature du triangle" selon Aaron Swartz, ce qui a motivé le développement de systèmes comme BitDNS puis Namecoin.}

\item \textbf{Question 92:} A
\\ \textit{Groupcoin/Devcoin se distinguait par son modèle de distribution particulier où 90\% des nouvelles unités créées étaient allouées aux développeurs impliqués dans le monde du logiciel libre, tandis que seulement 10\% étaient réservés aux mineurs du système. Bien que présenté comme une 'monnaie éthique', ce modèle constituait en réalité une aubaine pour les développeurs bénéficiaires, qui étaient choisis par le meneur du projet Unthinkingbit. Les autres options décrivent des modèles qui n'existaient pas à l'époque ou qui ne correspondent pas au fonctionnement réel de Groupcoin/Devcoin.}

\item \textbf{Question 93:} C
\\ \textit{Luke-Jr a effectivement mené une attaque de censure contre Coiledcoin le 6 janvier 2012 en utilisant sa coopérative de minage Eligius, ce qui s'est avéré fatal pour cette cryptomonnaie alternative. Il justifiait cette action par sa volonté de protéger l'écosystème Bitcoin contre les systèmes pyramidaux qui nuisent à sa réputation. L'action DMCA mentionnée dans une des mauvaises réponses a bien eu lieu, mais elle visait SolidCoin trois jours plus tard, pas Coiledcoin.}

\item \textbf{Question 94:} D
\\ \textit{BitDNS était spécifiquement conçu pour résoudre le problème de la centralisation des noms de domaine sous le contrôle de l'ICANN. Appamatto proposait d'utiliser une blockchain similaire à Bitcoin pour créer un système DNS décentralisé où les noms de domaine seraient gérés de manière distribuée, résistant ainsi aux saisures et à la censure. Cette proposition visait à appliquer le modèle décentralisé de Bitcoin au système de noms de domaine Internet, permettant de 'réaliser la quadrature du triangle de Zooko' selon Aaron Swartz.}

\item \textbf{Question 95:} D
\\ \textit{Tenebrix et Litecoin ont tous deux souffert du même problème technique : leur algorithme d'ajustement de difficulté était limité et ne pouvait multiplier ou diviser la difficulté que par 4. Cette limitation, combinée à une difficulté initiale trop basse, a causé un minage accéléré massif. Pour Tenebrix, 6 048 blocs ont été minés en 2 jours au lieu de 21, tandis que pour Litecoin, plus de 10 000 blocs ont été minés en moins de 24 heures au lieu de 12,5 jours prévus. Charlie Lee avait simplement repris le même paramètre défaillant de Tenebrix sans le corriger.}

\item \textbf{Question 96:} A
\\ \textit{Scrypt était une fonction de dérivation de clés conçue en 2009 par Colin Percival pour Tarsnap. Sa principale caractéristique était sa forte demande de mémoire, ce qui permettait de réduire drastiquement l'efficacité des GPU dans le minage et de rendre les CPU à nouveau profitables. Cette innovation technique visait à créer une cryptomonnaie 'favorable aux CPU et hostile aux GPU' selon les termes de Lolcust, contrairement aux autres réponses qui décrivent des caractéristiques inexistantes de Scrypt.}

\item \textbf{Question 97:} D
\\ \textit{Le texte indique clairement que Ross a remplacé le taux fixe de 6,23\% par un système dégressif allant de 15\% pour les petites ventes à 3,7\% pour les plus importantes. L'ancien système à taux fixe avantageait les petites ventes au détriment des grosses transactions, ce qui limitait la croissance du chiffre d'affaires de la plateforme. Cette modification visait à encourager les vendeurs à réaliser de plus gros volumes.}

\item \textbf{Question 98:} B
\\ \textit{LocalBitcoins a introduit en été 2012 un système de dépôt fiduciaire (escrow) innovant où les bitcoins sont conservés par la plateforme et peuvent être débloqués par SMS pour être versés à l'acheteur une fois que le vendeur confirme avoir reçu les espèces. Cette solution technique résolvait le problème majeur de confiance dans les échanges peer-to-peer en éliminant le risque que l'une des parties ne respecte pas sa part du marché. Les autres options proposées ne correspondent pas aux innovations réellement mises en place par Jeremias Kangas et son équipe.}

\item \textbf{Question 99:} D
\\ \textit{BitPay et BitInstant avaient des modèles économiques opposés. BitPay permettait aux commerçants d'accepter les bitcoins tout en étant payés en dollars, l'entreprise prenant le risque de volatilité en avançant les dollars et vendant les bitcoins plus tard. BitInstant facilitait plutôt l'achat de bitcoins via des dépôts d'espèces. LocalBitcoins gérait les échanges entre particuliers, et aucune de ces entreprises ne se concentrait sur les places de marché clandestines comme activité principale.}

\item \textbf{Question 100:} B
\\ \textit{Ross Ulbricht a choisi le nom Dread Pirate Roberts en référence au film The Princess Bride de 1987. Dans ce film, le nom n'est pas celui d'une personne mais un titre transmissible d'un pirate à l'autre, conçu pour 'inspirer la peur nécessaire'. Cette référence était stratégiquement pertinente car elle permettait de faciliter une transition entre différentes personnes ou de brouiller les pistes sur l'identité réelle de l'administrateur de Silk Road, une technique qui sera d'ailleurs utilisée plus tard.}

\item \textbf{Question 101:} D
\\ \textit{Selon l'étude de Chainalysis mentionnée dans le texte, Silk Road représentait entre 10 et 20 \% de l'activité économique totale sur la blockchain Bitcoin en 2012, avec des pics notables en février et septembre. Cette proportion était significative mais pas majoritaire, contrairement à ce que beaucoup pensaient à l'époque. La plateforme a ensuite vu sa part diminuer rapidement sous les 5 \% au profit d'autres activités comme le commerce conventionnel, le jeu d'argent et la spéculation. Les autres pourcentages proposés sont soit surestimés soit sous-estimés par rapport aux données réelles de cette analyse rétrospective.}

\item \textbf{Question 102:} C
\\ \textit{Le BIP 21 standardisait le format des URI bitcoin pour inclure toutes les informations nécessaires au paiement dans un seul code QR. Par exemple, 'bitcoin:175tWpb8K1S7NmH4Zx6rewF9WQrcZv245W?amount=20.3&label=Luke-Jr' contenait l'adresse de destination, le montant exact et une description. Cette standardisation était cruciale pour l'interopérabilité entre différents portefeuilles et facilitait grandement l'adoption commerciale en simplifiant le processus de paiement pour les consommateurs.}

\item \textbf{Question 103:} A
\\ \textit{L'étude de Nicolas Christin montre clairement que le volume journalier de Silk Road a doublé entre mars (8 000 BTC) et mai (15 000 BTC), puis a reculé à 11 000 BTC en juillet. Cette baisse apparente en bitcoins s'explique par la hausse du prix du bitcoin qui est passé de 5$ à plus de 8$ durant l'été, maintenant ainsi un chiffre d'affaires record de 1,9 million de dollars mensuels. Cette évolution démontre la sensibilité du volume transactionnel aux fluctuations de prix de la cryptomonnaie.}

\item \textbf{Question 104:} C
\\ \textit{Le texte explique clairement que le change entre bitcoin et argent liquide permet d'acheter ou de vendre du bitcoin 'de manière totalement confidentielle, sans laisser de trace de virement bancaire'. Cette confidentialité était cruciale pour BTCKing qui revendait des bitcoins aux acheteurs de Silk Road. Les autres options sont incorrectes : les frais de BitInstant étaient de 5\% (pas 2\%), il y avait une limite de 1000$ par dépôt (pas 50 000$), et aucune mention n'est faite de conversion en monnaies étrangères.}

\item \textbf{Question 105:} C
\\ \textit{Charlie Shrem justifiait sa collaboration avec BTCKing en affirmant que ce dernier n'avait enfreint aucune loi et que Silk Road elle-même n'était pas illégale, tout en soulignant les profits substantiels que cette relation apportait à BitInstant. Cette justification révèle l'ambiguïté morale et légale de l'époque, où les entreprises Bitcoin naviguaient dans une zone grise réglementaire. Les autres réponses sont incorrectes car BTCKing avait justement dépassé les limites de dépôt, son volume était très significatif (40 000$ par semaine), et il n'avait aucune licence officielle - il opérait comme revendeur sur Silk Road.}

\item \textbf{Question 106:} A
\\ \textit{L'étude de Nicolas Christin, qui révélait le succès économique impressionnant de Silk Road avec un chiffre d'affaires record de 1,9 million de dollars mensuels, a eu un effet paradoxal mais prévisible. Contrairement à ce qu'on pourrait attendre d'une exposition académique d'une plateforme illégale, cette étude a en fait servi de publicité gratuite pour Silk Road. Ross Ulbricht lui-même confirme cet impact en parlant de « croissance explosive » en novembre 2012, expliquant que la plateforme se trouvait « en terre inconnue en ce qui concerne le nombre d'utilisateurs ». Cette situation illustre comment la documentation scientifique d'un phénomène peut paradoxalement contribuer à son expansion.}

\item \textbf{Question 107:} A
\\ \textit{Le modèle économique de BitPay était révolutionnaire car il inversait le risque de volatilité. Contrairement à BitInstant qui aidait les gens à acheter des bitcoins, BitPay avançait immédiatement les dollars aux commerçants dès réception du paiement en bitcoin, puis se chargeait lui-même de vendre ces bitcoins sur le marché. Cette approche permettait aux commerçants de se prémunir complètement contre la volatilité du bitcoin tout en bénéficiant de ses avantages (pas de rétrofacturation, frais réduits, acceptation internationale).}

\item \textbf{Question 108:} B
\\ \textit{LocalBitcoins a intégré durant l'été 2012 un système de dépôt fiduciaire où les bitcoins sont conservés par la plateforme et peuvent être débloqués par SMS pour être versés à l'acheteur lorsque le vendeur confirme avoir reçu les espèces. Cette innovation était cruciale car elle résolvait le problème de confiance dans les échanges peer-to-peer en agissant comme tiers de confiance. Les autres options sont incorrectes : la plateforme ne demandait pas de vérification d'identité bancaire, n'utilisait pas de géolocalisation, et ne proposait pas d'assurance automatique.}

\item \textbf{Question 109:} D
\\ \textit{SatoshiDICE était révolutionnaire car il était « assurément équitable » (provably fair). Le système fonctionnait en publiant d'abord les empreintes (hash) des secrets du jour sur la blockchain, puis en dévoilant ces secrets quelques jours plus tard. Cela permettait aux joueurs de recalculer eux-mêmes les résultats en utilisant l'algorithme HMAC-SHA-512 avec le secret révélé et l'identifiant de leur transaction. Les autres options sont incorrectes : les adresses fixes ne garantissaient pas l'équité par elles-mêmes, il n'y avait pas de validation par nœuds multiples spécifique au jeu, et la simple transparence de la blockchain ne suffisait pas à prouver l'équité sans le mécanisme de secrets.}

\item \textbf{Question 110:} B
\\ \textit{Le 15 avril 2011, appelé « vendredi noir », marque la fermeture par le FBI des trois principales plateformes de poker aux États-Unis : PokerStars, Full Tilt Poker et Absolute Poker. Ces sites étaient accusés d'avoir contourné l'Unlawful Internet Gambling Enforcement Act en se faisant passer pour des commerçants vendant des marchandises comme des bijoux. Cette répression a créé un vide dans le marché du poker en ligne, ouvrant la voie à Bitcoin comme monnaie alternative pour contourner les restrictions bancaires traditionnelles.}

\item \textbf{Question 111:} A
\\ \textit{SatoshiDICE fonctionnait en envoyant systématiquement deux transactions pour chaque pari : une transaction entrante du joueur vers l'adresse du casino, puis une transaction sortante renvoyant soit les gains, soit une petite fraction confirmant le traitement. Cette mécanique fit exploser le nombre de transactions quotidiennes, passant de 7 000 à plus de 60 000 en juin 2012, représentant près de la moitié de toutes les transactions Bitcoin. Les autres options sont incorrectes car les secrets étaient stockés occasionnellement et non systématiquement, les calculs se faisaient hors chaîne, et les adresses fixes ne créaient pas de conflits de routage.}

\item \textbf{Question 112:} C
\\ \textit{Erik Voorhees a acquis le concept de SatoshiDICE pour 40 bitcoins, soit environ 200 $ au moment de l'échange. Les autres montants correspondent à différentes étapes : 45 BTC était l'investissement initial de Fireduck, 146 BTC étaient les profits générés en une semaine, et 34 500 BTC correspondent aux fonds levés lors de l'IPO informelle sur MPEx en août 2012, bien après l'acquisition initiale.}

\item \textbf{Question 113:} B
\\ \textit{Selon le chapitre, Bitcoin était particulièrement adapté au jeu d'argent car il résolvait les problèmes spécifiques rencontrés par ce secteur : d'une part, il réduisait les risques de rétrofacturation (chargebacks) qui posaient problème aux opérateurs, et d'autre part, il limitait la capacité des gouvernements à bloquer les paiements, comme cela s'était produit avec l'Unlawful Internet Gambling Enforcement Act aux États-Unis. Les autres options sont incorrectes car Bitcoin n'était pas sans frais, n'offrait qu'un anonymat relatif, et était lui-même sujet aux fluctuations de prix.}

\item \textbf{Question 114:} B
\\ \textit{Selon le chapitre, les trois principales plateformes de poker contournaient la loi en se faisant passer pour des commerçants qui vendaient « des marchandises telles que des bijoux et des balles de golf » aux joueurs, afin de ne pas attirer les soupçons de leurs banques. Cette stratégie de déguisement commercial leur permettait d'éviter que les virements soient identifiés comme des paiements destinés au jeu d'argent en ligne, ce qui était interdit par la loi américaine.}

\item \textbf{Question 115:} C
\\ \textit{Selon le chapitre, SatoshiDICE représentait près de la moitié des transactions réalisées sur la chaîne Bitcoin, mais seulement environ 5\% du volume économique. Cette disproportion s'explique par le fait que chaque pari générait deux transactions (une pour la mise, une pour le résultat), créant un grand nombre de petites transactions. Les autres réponses surestiment soit la proportion des transactions, soit confondent volume transactionnel et volume économique.}

\item \textbf{Question 116:} B
\\ \textit{SatoshiDICE utilisait un mécanisme spécifique pour garantir l'équité vérifiable : le « nombre chanceux » était obtenu en sélectionnant les deux premiers octets du condensé calculé par HMAC-SHA-512 d'une information secrète et de l'identifiant de la transaction entrante. Ce processus produisait un résultat entre 0 et 65 536. Les empreintes des secrets étaient publiées sur la blockchain et les secrets révélés ultérieurement, permettant aux joueurs de vérifier rétrospectivement que le casino n'avait pas triché. Cette méthode constituait l'innovation du système « assurément équitable » (provably fair).}

\item \textbf{Question 117:} B
\\ \textit{Bien que le casino ait connu un succès immédiat avec près de 5 000 BTC de paris en 10 jours, Fireduck a rapidement pris peur face aux implications légales de son activité. Il cherchait explicitement un repreneur qui pourrait assumer les risques juridiques tout en lui versant un salaire pour continuer le travail technique. Cette crainte des conséquences légales était justifiée dans un contexte où les autorités américaines venaient de fermer les principales plateformes de poker en ligne.}

\item \textbf{Question 118:} A
\\ \textit{Erik Voorhees a réalisé une IPO informelle en émettant 100 millions de parts S.DICE sur MPEx, vendant 10\% pour 34 500 bitcoins. Ces parts donnaient droit à des dividendes issus des revenus du site. Cette levée de fonds lui a permis de refaire le site à neuf en octobre 2012 et de développer des applications mobiles. Les autres options sont incorrectes : il n'y avait pas de transfert de propriété légale, la motivation n'était pas liée aux critiques techniques, et il ne s'agissait pas de créer une nouvelle cryptomonnaie mais bien d'utiliser Bitcoin.}

\item \textbf{Question 119:} B
\\ \textit{SatoshiDICE révolutionnait le secteur du jeu en ligne grâce à son système « provably fair » (assurément équitable). Contrairement aux casinos traditionnels où les joueurs devaient faire confiance aveuglément, SatoshiDICE permettait la vérification mathématique de l'équité. Le système publiait les empreintes des secrets sur la blockchain, puis révélait ces secrets ultérieurement, permettant aux joueurs de recalculer eux-mêmes les résultats en utilisant l'algorithme HMAC-SHA-512. Cette transparence cryptographique constituait une innovation majeure qui distinguait fondamentalement SatoshiDICE des plateformes centralisées non vérifiables de l'époque.}

\item \textbf{Question 120:} D
\\ \textit{SatoshiDICE a introduit le concept « provably fair » (assurément équitable) qui permettait aux joueurs de vérifier mathématiquement l'équité des tirages. Le système publiait les empreintes des secrets sur la blockchain, puis révélait ces secrets plus tard, permettant aux utilisateurs de recalculer eux-mêmes les résultats avec HMAC-SHA-512. Cette transparence cryptographique était révolutionnaire car elle éliminait la nécessité de faire confiance aveuglément au casino, contrairement aux plateformes centralisées traditionnelles qui n'étaient pas vérifiables.}

\item \textbf{Question 121:} B
\\ \textit{Bitcoin Savings & Trust utilisait un système de taux progressifs pour encourager les gros dépôts : 100 BTC rapportaient 4,2\% par semaine tandis que 1000 BTC et plus rapportaient 7\% par semaine. Ce système pyramidal visait à attirer les investisseurs importants. Le taux de 1\% par jour était proposé au début, mais fut modifié. Les taux n'étaient ni dégressifs ni basés sur les fluctuations du bitcoin, mais fixes selon les paliers d'investissement.}

\item \textbf{Question 122:} A
\\ \textit{Le troisième piratage de Bitcoinica résultait directement de la publication du code source par Amir Taaki le 7 juillet 2012. Ce code contenait une clé API qui s'avérait également être le mot de passe du coffre-fort LastPass de Bitcoinica, donnant accès aux mots de passe de la plateforme, notamment celui de Mt. Gox. Les autres réponses décrivent des méthodes utilisées lors d'autres piratages : Rackspace pour le deuxième piratage de mai 2012, Linode pour le premier piratage de mars 2012.}

\item \textbf{Question 123:} D
\\ \textit{L'absence de sauvegarde de la base de données lors du piratage du 11 mai 2012 a créé une situation catastrophique où Bitcoinica ne pouvait plus identifier précisément les soldes de ses clients. Pour résoudre ce problème, la Bitcoin Consultancy a mis en place un système complexe de demandes de remboursement où chaque créancier devait fournir des preuves de ses dépôts : copies d'emails, codes Mt. Gox utilisés, identifiants de transactions, pièces d'identité, etc. Le plan prévoyait de payer 50\% des fonds initialement, puis le reste plus tard pour limiter les erreurs. Cette approche était la seule solution viable face à l'absence totale de données fiables sur les comptes clients.}

\item \textbf{Question 124:} B
\\ \textit{Le piratage de Linode s'est produit par l'accès au portail du service client, permettant au pirate de prendre le contrôle de huit comptes liés à Bitcoin. Cette méthode diffère des autres piratages de l'époque : contrairement au troisième piratage de Bitcoinica qui exploitait une clé API dans le code source, ou aux infections par logiciels malveillants comme celle d'Allinvain, le pirate de Linode a ciblé directement l'infrastructure de l'hébergeur pour accéder aux serveurs des clients.}

\item \textbf{Question 125:} B
\\ \textit{Zhou Tong a explicitement déclaré dans son message du 13 mai 2012 que sa décision de quitter Bitcoin n'était pas liée aux piratages de Bitcoinica, mais reposait sur une réflexion philosophique. Il considérait que la spéculation était un jeu à somme nulle et qu'il n'avait pas réussi à créer de la valeur pour la société. Il voulait désormais construire des produits qui économisent du temps et de l'argent aux gens, estimant que Bitcoin ne répondait pas à ses besoins réels. Bien qu'il ait déménagé en Australie pour continuer ses études, cela n'était pas la raison principale de son départ de l'écosystème Bitcoin.}

\item \textbf{Question 126:} B
\\ \textit{L'escroquerie de Tony76 était unique car elle reposait entièrement sur l'ingénierie sociale plutôt que sur des failles techniques. Tony76 était un vendeur respecté avec plus de 500 ventes et une réputation excellente, ce qui lui permettait d'exiger que les acheteurs utilisent la finalisation anticipée (bypassing l'escrow). Cette méthode différait des autres fraudes de 2012 qui exploitaient généralement des vulnérabilités informatiques comme les piratages de serveurs ou les failles de sécurité.}

\item \textbf{Question 127:} D
\\ \textit{La fermeture brutale de Bitcoin Savings & Trust, qui gérait plus de 500 000 bitcoins (environ 5\% de tous les bitcoins en circulation), a provoqué une panique sur le marché. Les investisseurs ont massivement vendu leurs bitcoins, causant une chute vertigineuse du prix de 15,4$ à moins de 8$ en seulement 48 heures sur Mt. Gox. Cette réaction violente s'explique par l'ampleur du système pyramidal et la perte de confiance généralisée qui a suivi son effondrement.}

\item \textbf{Question 128:} A
\\ \textit{Le piratage de Bitfloor du 3-4 septembre 2012 a été rendu possible par le fait que les clés privées étaient stockées dans une sauvegarde non chiffrée sur les serveurs. Quand le pirate a accédé aux serveurs, il a pu directement exploiter cette sauvegarde pour voler plus de 24 000 BTC. Cette négligence fondamentale en matière de sécurité contraste avec les autres réponses qui, bien que représentant des vulnérabilités potentielles, n'étaient pas la cause directe de ce vol massif.}

\item \textbf{Question 129:} D
\\ \textit{Selon le texte, Trendon Shavers expliquait que son revenu provenait de trois sources : la revente locale, l'arbitrage de marché, et le prêt de bitcoins. Cependant, c'est cette dernière activité qui est rapidement devenue la principale source de revenus de son système. Les autres options, bien que mentionnées par Shavers comme sources de profit, n'étaient pas l'activité principale. L'arbitrage était même qualifié de spéculation et manipulation des cours plutôt que de véritable arbitrage.}

\item \textbf{Question 130:} D
\\ \textit{À son sommet, Bitcoin Savings & Trust gérait plus de 500 000 bitcoins sur un total de 9,7 millions de bitcoins en circulation, soit environ 5\% de l'offre totale. Cette proportion énorme démontre l'ampleur exceptionnelle de cette pyramide de Ponzi et explique pourquoi sa fermeture a eu un impact si dramatique sur le marché, provoquant une chute du prix de 15,4$ à moins de 8$ en 48 heures.}

\item \textbf{Question 131:} B
\\ \textit{Le troisième piratage de Bitcoinica résulte directement de la publication du code source par Amir Taaki le 7 juillet 2012. Ce code contenait une clé API qui s'avérait également être le mot de passe du coffre-fort LastPass de Bitcoinica, donnant accès aux mots de passe de la plateforme, notamment celui de Mt. Gox. Cette vulnérabilité a été exploitée le 13 juillet par un pirate qui a retiré 40 000 bitcoins et 40 000 LR-USD. Les autres options correspondent à des incidents antérieurs différents qui n'ont pas causé ce troisième piratage spécifique.}

\item \textbf{Question 132:} B
\\ \textit{La multiplication des vols et escroqueries a créé une image désastreuse de Bitcoin qui a eu un effet repoussoir sur la population générale et encouragé l'intervention étatique. Pour remédier à ces risques, il fallait faire un effort de communication et redorer le blason de Bitcoin. Cela s'est traduit par la création de la Fondation Bitcoin le 27 septembre 2012, date pivot de l'histoire de Bitcoin, dans le but de courtiser les autorités par le lobbying et rendre Bitcoin acceptable aux yeux du grand public.}


\end{enumerate}

\end{document}
